\subsection{半单模中的标量扩张}

命题 8. —— 设 A 是一个 K-代数，M 是一个 A-模，L 是域 K 的一个扩张。用 D 表示 M 的交换子，用 Z 表示 D 的中心。

\begin{enumerate}
\def\labelenumi{\alph{enumi})}
\item
  假设 \(A_{(L)}\) -模 \(M_{(L)}\) 是单的（相应地，同型的、半单的）。则 A-模 M 是单的（相应地，同型的、半单的）。
\item
  假设 A-模 M 是半单的，且 M 或 L 在 K 上是有限维的。\(\mathrm{A}_{(\mathrm{L})}\) -模 \(\mathrm{M}_{(\mathrm{L})}\) 是半单的当且仅当环 \(\mathrm{Z}_{(\mathrm{L})}\) 是约化的。\(\mathrm{A}_{(\mathrm{L})}\) -模 \(\mathrm{M}_{(\mathrm{L})}\) 是同型且非零的当且仅当环 \(\mathrm{Z}_{(\mathrm{L})}\) 是一个整环。
\item
  假设 A-模 M 是单的。\(\mathrm{A}_{(\mathrm{L})}\) -模 \(\mathrm{M}_{(\mathrm{L})}\) 是半单的（相应地，同型的、单的）当且仅当环 \(\mathrm{D}_{(\mathrm{L})}\) 是半单的（相应地，单的、一个体）。
\end{enumerate}

断言 a) 是命题 1（VIII, 第 211 页）的特殊情形，断言 b) 是命题 5（VIII, 第 219 页）的特殊情形，而断言 c) 是 VIII, 第 215 页的推论 1 的特殊情形。

推论 1. —— a) 假设 A-模 M 是半单的，K 的扩张 L 是可分的，且 M 或 L 在 K 上是有限维的。则 \(\mathrm{A}_{(\mathrm{L})}\) -模 \(\mathrm{M}_{(\mathrm{L})}\) 是半单的。

\begin{enumerate}
\def\labelenumi{\alph{enumi})}
\setcounter{enumi}{1}
\tightlist
\item
  假设 A-模 M 是单的且其交换子等于 K。则 \(\mathrm{A}_{(\mathrm{L})}\) -模 \(\mathrm{M}_{(\mathrm{L})}\) 是单的。
\end{enumerate}

断言 a) 由 VIII, 第 221 页的推论得出。断言 b) 是命题 8, c) 的特殊情形。

推论 2. —— 设 L 是域 K 的一个扩张。用 Z 表示 K-代数 A 的中心。

\begin{enumerate}
\def\labelenumi{\alph{enumi})}
\item
  若 L-代数 \(A_{(L)}\) 是半单的，则 K-代数 A 是半单的。
\item
  假设 K-代数 A 是半单的，且 L 或 A 在 K 上是有限维的。L-代数 \(\mathrm{A}_{(\mathrm{L})}\) 是半单的当且仅当环 \(\mathrm{Z}_{(\mathrm{L})}\) 是约化的；特别地，当 L 是 K 的一个可分扩张时即是这种情形。L-代数 \(\mathrm{A}_{(\mathrm{L})}\) 是单的当且仅当环 \(\mathrm{Z}_{(\mathrm{L})}\) 是一个整环；特别地，当 A 的中心等于 K 时即是这种情形。
\end{enumerate}

断言 a) 与 b) 由应用于 A-模 \(\mathrm{A}_s\) 的命题 8, a) 与 b) 得出。

命题 9. —— 设 A 是一个 K-代数，L 是 K 的一个可分扩张。

\begin{enumerate}
\def\labelenumi{\alph{enumi})}
\item
  若 M 是一个无根的 A-模，则 \(\mathrm{A}_{(\mathrm{L})}\) -模 \(\mathrm{M}_{(\mathrm{L})}\) 是无根的。
\item
  若 K-代数 A 是无根的，则 L-代数 \(\mathrm{A}_{(\mathrm{L})}\) 是无根的。
\end{enumerate}

我们来证明断言 a)。设 M 是一个无根的 A-模。我们把 M 与它在 \(\mathrm{M}_{(\mathrm{L})}\) 中的典范像等同起来。设 N 是 M 的一个极大子模。由于 A-模 M/N 是单的，由 VIII, 第 218 页的命题 4 得出，\(\mathrm{A}_{(\mathrm{L})}\) -模 \((\mathrm{M}/\mathrm{N})_{(\mathrm{L})} = \mathrm{M}_{(\mathrm{L})}/\mathrm{N}_{(\mathrm{L})}\) 是无根的，因此由 VIII, 第 152 页的推论 1, c) 有 \(\mathfrak{R}_{\mathrm{A}_{(\mathrm{L})}}(\mathrm{M}_{(\mathrm{L})}) \subset \mathrm{N}_{(\mathrm{L})}\)。另一方面，由 II, §7, No.~7, 第 306 页的命题 14 的推论得出，当 N 遍历 M 的极大子模所成之集时，诸 \(\mathrm{N}_{(\mathrm{L})}\) 的交化为 0。因此，\(\mathrm{A}_{(\mathrm{L})}\) -模 \(\mathrm{M}_{(\mathrm{L})}\) 是无根的。

断言 b) 由应用于 A-模 \(\mathrm{A}_s\) 的断言 a) 得出。

命题 10. —— 设 A 是一个 K-代数，L 是 K 的一个扩张。设 M 是一个 A-模。

\begin{enumerate}
\def\labelenumi{\alph{enumi})}
\tightlist
\item
  我们作以下两个假设之一：
\end{enumerate}

\begin{enumerate}
\def\labelenumi{(\roman{enumi})}
\item
  A-模 M 是有限生成的，且 L 在 K 上是代数的。
\item
  环 A 是左阿廷的。
\end{enumerate}

则我们有包含关系

\[
\mathfrak {R} _ {\mathrm{A}} (\mathrm{M}) _ {(\mathrm{L})} \subset \mathfrak {R} _ {\mathrm{A} _ {(\mathrm{L})}} (\mathrm{M} _ {(\mathrm{L})}).
\]

\begin{enumerate}
\def\labelenumi{\alph{enumi})}
\setcounter{enumi}{1}
\tightlist
\item
  若 L 是 K 的一个可分扩张，则我们有包含关系
\end{enumerate}

\[
\mathfrak {R} _ {\mathrm{A} _ {(\mathrm{L})}} (\mathrm{M} _ {(\mathrm{L})}) \subset \mathfrak {R} _ {\mathrm{A}} (\mathrm{M}) _ {(\mathrm{L})}.
\]

我们来证明断言 a)。考虑情形 (i)。首先假设 L 在 K 上次数有限。则 A-模 \(M_{(L)}\) 是有限生成的。用 f 表示从 A 到 \(A_{(L)}\) 的典范同态；环 \(A_{(L)}\) 由其中心与 \(f(A)\) 的并生成。因此我们可以把 VIII, 第 174 页的命题 3 应用于 \(A_{(L)}\) -模 \(M_{(L)}\)。这给出包含关系 \(\mathfrak{R}_{\mathrm{A}}(\mathrm{M}_{(\mathrm{L})}) \subset \mathfrak{R}_{\mathrm{A}_{(\mathrm{L})}}(\mathrm{M}_{(\mathrm{L})})\)，更有 \(\mathfrak{R}_{\mathrm{A}}(\mathrm{M})_{(\mathrm{L})} \subset \mathfrak{R}_{\mathrm{A}_{(\mathrm{L})}}(\mathrm{M}_{(\mathrm{L})})\)（VIII, 第 152 页，推论 1）。

现在来处理一般情形。设 \(x_{1},\ldots,x_{n}\) 是生成 A-模 M 的元素，x 是 \(\mathfrak{R}_{\mathrm{A}}(\mathrm{M})\) 的一个元素。设 \(a_{1},\ldots,a_{n}\) 是 \(\mathrm{A}_{(\mathrm{L})}\) 的元素；由于 L 在 K 上是代数的，存在 K 的一个包含在 L 中的有限扩张 \(L'\)，使得诸 \(a_{i}\) 属于 \(\mathrm{A}_{(\mathrm{L}')}\)。由前所述，x 属于 \(\mathrm{A}_{(\mathrm{L}')}\) -模 \(\mathrm{M}_{(\mathrm{L}')}\) 的根；由 VIII, 第 153 页的推论得出，诸元素 \(x_{i}+a_{i}x\)（\(1\leqslant i\leqslant n\)）生成 \(\mathrm{A}_{(\mathrm{L}')}\) -模 \(\mathrm{M}_{(\mathrm{L}')}\)，从而生成 \(\mathrm{A}_{(\mathrm{L})}\) -模 \(\mathrm{M}_{(\mathrm{L})}\)。由同一推论，x 属于 \(\mathrm{A}_{(\mathrm{L})}\) -模 \(\mathrm{M}_{(\mathrm{L})}\) 的根。

现在考虑情形 (ii)。设 \(\mathfrak{r}\) 是 A 的根，于是 A-模 M 的根等于 \(\mathfrak{rM}\)（VIII, 第 174 页，推论）。A 的根 \(\mathfrak{r}\) 是 A 的一个幂零双边理想（VIII, 第 173 页，命题 1）；因此 \(\mathfrak{r}_{(L)}\) 是 \(A_{(L)}\) 的一个幂零双边理想。由此得出 \(\mathfrak{r}_{(L)} \subset \Re(A_{(L)})\)（VIII, 第 156 页，定理 1），且 VIII, 第 158 页的命题 6 蕴含 \(\Re(A_{(L)})M_{(L)} \subset \Re_{A_{(L)}}(M_{(L)})\)。于是我们有 \(\Re_A(M) = \mathfrak{rM} \subset \Re_{A_{(L)}}(M_{(L)})\)，这就完成了断言 a) 的证明。

模 \(\mathrm{M}/\mathfrak{R}_{\mathrm{A}}(\mathrm{M})\) 是无根的。若 L 是 K 的一个可分扩张，则由 VIII, 第 223 页的命题 9 得出，\(\mathrm{A}_{(\mathrm{L})}\) -模 \((\mathrm{M}/\mathfrak{R}_{\mathrm{A}}(\mathrm{M}))_{(\mathrm{L})}\) 是无根的。因此，我们有包含关系

\[
\mathfrak {R} _ {\mathrm{A} _ {(\mathrm{L})}} (\mathrm{M} _ {(\mathrm{L})}) \subset \mathfrak {R} _ {\mathrm{A}} (\mathrm{M}) _ {(\mathrm{L})}.
\]

推论. —— 设 L 是 K 的一个可分扩张。若 L 在 K 上是代数的，或者环 A 是左阿廷的，则我们有 \(\mathfrak{R}(A_{(L)}) = \mathfrak{R}(A)_{(L)}\)。

这是命题 10 中 \(\mathrm{M} = \mathrm{A}_s\) 的情形。

\BourbakiExercises{习题}

\begin{enumerate}
\def\labelenumi{\arabic{enumi})}
\tightlist
\item
  设 A 与 B 是 K-代数，M 是一个 A-模；把 \(M \otimes B\) 视为代数 \(C = A \otimes B\) 上的模。把 M 等同为 \(M \otimes B\) 的一个 A-子模。
\end{enumerate}

\begin{enumerate}
\def\labelenumi{\alph{enumi})}
\item
  证明包含关系 \(M \cap \mathfrak{R}_{\mathrm{C}}(\mathrm{M} \otimes \mathrm{B}) \subset \mathfrak{R}_{\mathrm{A}}(\mathrm{M})\)（考虑一个从 M 到某个单 A-模的 A-线性同态）。
\item
  证明在下列每种情形中等号成立：
\end{enumerate}

\begin{enumerate}
\def\labelenumi{(\roman{enumi})}
\item
  A 在 \(\mathbf{K}\) 上是有限维的。
\item
  A-模 M 是有限生成的，且 A 是一个由在 K 上有限维的子代数组成的右向定向族的并。
\item
  \(M = A_{s}\)，且 A 的根是幂零的。
\end{enumerate}

（在情形 (i) 中，可证明每个从 \(\mathrm{M} \otimes \mathrm{B}\) 到单 C-模 S 的同态都在 \(\Re_{\mathrm{A}}(\mathrm{M})\) 上为零，办法是注意到 \(\Re_{\mathrm{C}}(\mathrm{S})\) 为零。）

\begin{enumerate}
\def\labelenumi{\arabic{enumi})}
\setcounter{enumi}{1}
\item
  设 \(\mathrm{A}\) 是一个为整环的 K-代数，\(\mathrm{L}\) 是它的分式域。证明 \(\mathrm{L}\) 是一个单 \(\mathrm{A} \otimes \mathrm{L}\) -模，但它作为 A-模的根等于 \(\mathrm{L}\)。
\item
  给出（非零）K-代数 \(\mathrm{A}_1\) 与 \(\mathrm{A}_2\) 的一个例子，使得 \(\Re (\mathrm{A}_1)\neq 0\) 而 \(\Re (\mathrm{A}_1\otimes \mathrm{A}_2) = 0\)（参见 VIII, 第 167 页，习题 11）。
\item
  给出一个交换域 K 与 K 的两个扩张 E 和 F 的例子，使得环 \(A = E \otimes_{K} F\) 是半单的但不是域（相应地，不是半单的）。推出 A-模 \(A_{s} = E_{s} \otimes_{K} F_{s}\) 是半单的但不是单的（相应地，不是半单的）。
\end{enumerate}

¶ 5) 设 A 是一个拟单 K-代数（VIII, 第 120 页，注 2）。

\begin{enumerate}
\def\labelenumi{\alph{enumi})}
\item
  证明 (A, A)-双模 \(_{s}A_{d}\) 是单的；推出 A 的中心 Z 是一个体，同构于该双模的交换子。
\item
  设 B 是一个 K-代数。证明 \(A \otimes B\) 的每个双边理想都具有形式 \(A \otimes_{Z} b\)，其中 b 是 B 的一个双边理想（化归到情形 Z = K，并利用 VIII, 第 63 页的推论 2）。
\item
  推出若 A 与 B 是拟单 K-代数且其中之一的中心为 K，则代数 \(\mathrm{A} \otimes \mathrm{B}\) 是拟单的。
\end{enumerate}

\(d\) ) 设 B 是一个 K-代数，A 是 B 的一个中心 Z 包含 K 的拟单子代数。证明 A 与它在 B 中的交换子 \(\mathrm{A}'\) 在 Z 上线性分离（利用 \(b\)）。

\begin{enumerate}
\def\labelenumi{\alph{enumi})}
\setcounter{enumi}{4}
\item
  设 A 是一个拟单 K-代数，Z 是它的中心。设 \(\varphi : \mathrm{A} \otimes \mathrm{A}^{\circ} \to \operatorname{End}_{\mathrm{Z}}(\mathrm{A})\) 是把 \(a \otimes b\) 映为自同态 \(x \mapsto axb\) 的同态。证明 \(\varphi\) 是单射且它的像 H 是一个拟单代数（利用 \(c\) 与 \(d\)）。此外证明 H 在 \(\operatorname{End}_{\mathrm{Z}}(\mathrm{A})\) 中稠密（VIII, 第 129 页，习题 3）；我们有 \(\mathrm{H} = \operatorname{End}_{\mathrm{Z}}(\mathrm{A})\) 当且仅当 A 在 Z 上是有限维的（参见 VIII, 第 128 页，习题 2）。
\item
  设 A 是一个以 K 为中心的半单代数。由 e) 推出，K-代数 \(A \otimes A^{o}\) 是半单的当且仅当 A 在 K 上秩有限（参见 VIII, 第 238 页，定理 2）。
\end{enumerate}

\begin{enumerate}
\def\labelenumi{\arabic{enumi})}
\setcounter{enumi}{5}
\item
  设 A 是一个具有非零幂零根的 K-代数。证明同态 \(\varphi: A \otimes A^{\circ} \to \operatorname{End}_{K}(A)\)（习题 5）不是单射。
\item
  设 \(A_{1}\) 与 \(A_{2}\) 是半单 K-代数，\(Z_{1}\) 与 \(Z_{2}\) 是它们的中心。假设 \(A_{1}\) 在 K 上次数有限，且环 \(Z_{1} \otimes Z_{2}\) 是约化的。设 f 是一个从 \(A_{2}\) 到 \(A_{1}\) 的 K-代数同态。证明 \(f(\mathrm{A}_{2})\) 在 \(A_{1}\) 中的交换子是一个半单 K-代数（把 \(A_{1}\) 视为 \((\mathrm{A}_{2} \otimes \mathrm{A}_{1}^{\circ})\) -模，并利用 VIII, 第 221 页的命题 7）。
\item
  设 K 是一个交换域，L 是 K 的一个次数 \textgreater{} 1 的有限根式扩张。设 V 是 L 上的一个无限维向量空间，A 是 \(\operatorname{End}_{\mathrm{L}}(\mathrm{V})\) 的由有限秩自同态与恒等映射生成的 K-子代数。证明 L-代数 \(\mathrm{A}_{(\mathrm{L})}\) 具有非零根，尽管其中心等于 L（利用习题 7）。
\item
  设 A 与 B 是 K-代数，M 是一个自由 A-模，N 是一个 B-模。
\end{enumerate}

\begin{enumerate}
\def\labelenumi{\alph{enumi})}
\item
  证明典范同态 \(\vartheta\)（从 \(\mathrm{End}_{\mathrm{A}}(\mathrm{M})\otimes \mathrm{End}_{\mathrm{B}}(\mathrm{N})\) 到 \(\mathrm{End}_{\mathrm{A}\otimes \mathrm{B}}(\mathrm{M}\otimes \mathrm{N})\)）的像是环 \(\mathrm{End}_{\mathrm{A}\otimes \mathrm{B}}(\mathrm{M}\otimes \mathrm{N})\) 的一个稠密子环。
\item
  若 B-模 N 的位似环在其双交换子中稠密，则 \((A \otimes B)\) -模 \(M \otimes N\) 的位似环亦然。
\item
  假设存在 N 的一个元素 y 与 N 的自同态的一个序列 \((v_{n})\)，使得诸元素 \(v_{n}(y)\) 生成 K 上的一个无限维向量空间。证明映射 \(\vartheta\) 不是满射。
\end{enumerate}

\begin{enumerate}
\def\labelenumi{\arabic{enumi})}
\setcounter{enumi}{9}
\tightlist
\item
  保留习题 9 的假设，并此外假设 B-模 N 是单的；用 D 表示它的交换子。
\end{enumerate}

\begin{enumerate}
\def\labelenumi{\alph{enumi})}
\item
  映射 \(\vartheta\) 是满射的当且仅当 M 在 A 上具有有限基，或者 D 在 K 上是有限维的（利用习题 9, c)）。
\item
  证明 \((\mathrm{A} \otimes \mathrm{B})\) -模 \(M \otimes N\) 是半单的当且仅当环 \(A \otimes D^{o}\) 是半单的（化归到情形 \(M = A_{s}\)；注意到 \(M \otimes N\) 是一个自由 \(A \otimes D^{o}\) -模，并利用 VIII, 第 85 页的命题 5 以及习题 9, a)）。
\end{enumerate}

¶ 11) 设 A 与 B 是 K-代数；假设环 A 与 B 都是本原的，且它们各自的底座 s 与 t（VIII, 第 130 页，习题 8）不化为 0。设 D（相应地，E）是 A（相应地，B）的一个极小左理想 a（相应地，b）的交换子；假设代数 D ⊗ E 是单的。证明 A ⊗ B 是一个以 s ⊗ t 为底座的本原环（注意到 (A ⊗ B)-模 a ⊗ b 是有限长度的同型模；利用 VIII, 第 131 页的习题 9, a)）。

\begin{enumerate}
\def\labelenumi{\arabic{enumi})}
\setcounter{enumi}{11}
\tightlist
\item
  设 \(A_{1}\) 是 K-代数 \(\mathrm{K}(\mathrm{X})[\mathrm{U}]_{1,d/d\mathrm{X}}\)（VIII, 第 11 页），\(A_{2}\) 是 K-代数 \(K[[T]]\)，A 是 K-代数 \(A_{1} \otimes_{K} A_{2}\)。把 \(A_{1}\) 与 \(A_{2}\) 等同为 A 的子代数。
\end{enumerate}

\begin{enumerate}
\def\labelenumi{\alph{enumi})}
\item
  证明 \(A_{2}\) 是 A 的中心，A 不含任何零因子，且 A 的可逆元恰是 \(A_{2}\) 的那些元素。
\item
  证明 A 的双边理想恰是诸理想 \(AT^{n}\)（\(n \geqslant 0\)）。
\item
  证明环 A 是本原的，尽管其中心具有非零根（注意到 A 的一个包含 1 - TU 的极大左理想不含 0 以外的任何双边理想）。
\end{enumerate}

\begin{enumerate}
\def\labelenumi{\arabic{enumi})}
\setcounter{enumi}{12}
\tightlist
\item
  设 D 与 E 是中心包含 K 的体，V（相应地，W）是 D（相应地，E）上的一个向量空间。证明 \((\mathrm{D} \otimes_{\mathrm{K}} \mathrm{E})\) -模 \(V \otimes_{K} W\) 是自由的。
\end{enumerate}

¶ 14) 设 V 是一个体 D 上的向量空间。用 \(\Omega\) 表示加群 V 的自同态环，并把 D 等同为 \(\Omega\) 的一个子环。令 \(A = \text{End}_{\text{D}}(\text{V})\)。

\begin{enumerate}
\def\labelenumi{\alph{enumi})}
\item
  证明 A 与 D 在它们的公共中心 Z 上线性分离（参见习题 5, d)）。
\item
  \(D \otimes_{Z} A\) 在 \(\Omega\) 中的像等于 \(\operatorname{End}_{\mathbb{Z}}(\mathrm{V})\) 当且仅当 D 在 Z 上秩有限。（为看出该条件是必要的，首先由习题 5, e) 推出 V = D 的情形，然后固定 V 的一个元素 \(x \neq 0\)，并考虑 \(\operatorname{End}_{\mathbb{Z}}(\mathrm{V})\) 的子代数 B，它由满足 \(u(\mathrm{D}x) \subset \mathrm{D}x\) 的 V 的 Z-自同态 u 组成。注意到等式 \(\operatorname{End}_{\mathbb{Z}}(\mathrm{V}) = \mathrm{DA}\) 蕴含 \(\mathrm{B} = \mathrm{D}(\mathrm{B} \cap \mathrm{A})\)。）
\end{enumerate}

\begin{enumerate}
\def\labelenumi{\arabic{enumi})}
\setcounter{enumi}{14}
\tightlist
\item
  \begin{enumerate}
  \def\labelenumii{\alph{enumii})}
  \tightlist
  \item
    设 \(\mathrm{K}\) 是一个交换域，A 是一个 K-代数，\(\mathrm{B} = \mathrm{K}[\mathrm{X}_{\iota}]_{\iota \in \mathrm{I}}\) 是 \(\mathrm{K}\) 上的一个多项式代数。证明若 \(\mathrm{A}\) 的每个非零元素都是正则的，则代数 \(\mathrm{A} \otimes_{\mathrm{K}} \mathrm{B}\) 具有同样的性质。
  \end{enumerate}
\end{enumerate}

\begin{enumerate}
\def\labelenumi{\alph{enumi})}
\setcounter{enumi}{1}
\tightlist
\item
  推出若 D 是一个在 K 上秩有限的体，而 E 是 K 的一个纯超越扩张，则环 \(D \otimes_{K} E\) 是一个体。
\end{enumerate}

¶ 16) 设 D 是一个体，V 是 D 上的一个无限维右向量空间，A 是环 \(\mathrm{End}_{\mathrm{D}}(\mathrm{V})\) 的一个底座 \(\mathfrak{s}\) 非零的稠密子环（VIII, 第 130 页，习题 8 与 9）。设 K 是 D 的中心的一个子域，B 是一个 K-代数。证明若 K-代数 \(\mathrm{D}^{\circ} \otimes \mathrm{B}\) 的根非零，则 \(\mathrm{A} \otimes \mathrm{B}\) 的根亦非零。（设 \(\mathfrak{a}\) 为元素 \(\sum u_{i} \otimes b_{i}\)（属于 \(\mathfrak{s} \otimes \mathrm{B}\)）中满足 \(\sum \langle u_{i}(x), x^{*} \rangle \otimes b_{i} \in \Re(\mathrm{D}^{\circ} \otimes \mathrm{B})\)（对所有 \(x \in \mathrm{V}\) 与 \(x^{*} \in \mathrm{V}^{*}\)）者所成之集。证明 \(\mathfrak{a}\) 是 \(\mathrm{A} \otimes \mathrm{B}\) 的一个双边理想，并利用 VIII, 第 167 页的习题 10, e) 与 VIII, 第 131 页的习题 9，证明它非零且包含在 \(\mathrm{A} \otimes \mathrm{B}\) 的根中。）

\begin{enumerate}
\def\labelenumi{\arabic{enumi})}
\setcounter{enumi}{16}
\tightlist
\item
  设 K 是一个交换域，D 与 E 是中心包含 K 的体。假设 \([E:K]=m\) 有限，且体 D 与 E 之一的中心等于 K。则 K-代数 \(D\otimes_{K}E\) 同构于一个体 C 上的有限维向量空间 h 的自同态代数（VIII, 第 222 页，推论 2）。
\end{enumerate}

\begin{enumerate}
\def\labelenumi{\alph{enumi})}
\item
  证明 \(h\) 整除 \(m\)（参见 VIII, 第 205 页，公式 (32)）。
\item
  设 W 是 D 上的一个向量空间；则 E 同构于 \(\operatorname{End}_{\mathrm{D}}(\mathrm{W})\) 的一个 K-子代数当且仅当 W 的维数是无限的或是 m/h 的倍数（注意到此时 W 是一个 \((\mathrm{D} \otimes_{\mathrm{K}} \mathrm{E})\) -模）。
\item
  证明若 \([\mathrm{D}:\mathrm{K}]\) 有限且与 \(m\) 互素，则 \(\mathrm{D}\otimes_{\mathrm{K}}\mathrm{E}\) 是一个体（应用 \(a\)）。
\end{enumerate}

\section{绝对半单代数}

\subsection{绝对半单模}

定义 1. —— 设 K 是一个交换域，A 是一个 K-代数。若对 K 的每个扩张 L，\(\mathrm{A}_{(\mathrm{L})}\) -模 \(\mathrm{M}_{(\mathrm{L})}\) 都是半单的，则称 A-模 M 是绝对半单的。

每个绝对半单模都是半单的。反之，若域 K 是完满的，则每个在 K 上有限维的半单 A-模都是绝对半单的（VIII, 第 222 页，推论 1, a)），因为完满域的每个扩张都是可分的（V, §7, No.~1, 第 36 页，命题 2）。

命题 1. —— 设 K 是一个交换域，A 是一个 K-代数。

\begin{enumerate}
\def\labelenumi{\alph{enumi})}
\item
  绝对半单 A-模的任意直和是一个绝对半单 A-模。绝对半单模的每个子模与商模都是绝对半单的。
\item
  设 M 是一个 A-模，L 是域 K 的一个扩张。A-模 M 是绝对半单的当且仅当 \(\mathrm{A}_{(\mathrm{L})}\) -模 \(\mathrm{M}_{(\mathrm{L})}\) 是绝对半单的。
\end{enumerate}

断言 a) 由关于半单模的类似断言得出（VIII, 第 56 页，推论 1 与 3）。

假设 A-模 M 是绝对半单的，并设 \(L'\) 是 L 的一个扩张。作为 \(\mathrm{A}_{(\mathrm{L}')}\) -模，\(L' \otimes_{L} M_{(L)}\) 同构于 \(\mathrm{M}_{(\mathrm{L}')}\)（II, §5, No.~1, 第 273 页，命题 2）；因此它是一个半单模。这就证明了 \(\mathrm{M}_{(\mathrm{L})}\) 是绝对半单的。

反之，假设 \(M_{(L)}\) 是绝对半单的。设 \(L'\) 是 K 的一个扩张。存在一个复合扩张 \((\Omega, u, v)\)（L 与 \(L'\) 的；V, §2, No.~4, 第 13 页，推论）；我们把 L 与 \(L'\) 等同为 \(\Omega\) 的子扩张。

\(A_{(\Omega)}\) -模 \(M_{(\Omega)}\) 同构于 \((\mathrm{M}_{(\mathrm{L})})_{(\Omega)}\)；因此它是半单的。但 \(M_{(\Omega)}\) 也同构于 \((\mathrm{M}_{(\mathrm{L}^{\prime})})_{(\Omega)}\)，且 VIII, 第 222 页的命题 8, a) 蕴含 \(M_{(L^{\prime})}\) 是半单的。于是 M 是绝对半单的。

命题 2. —— 设 K 是一个交换域，A 是一个 K-代数，M 是一个在 K 上有限维的 A-模。下列性质等价：

\begin{enumerate}
\def\labelenumi{(\roman{enumi})}
\item
  A-模 M 是绝对半单的。
\item
  存在 \(\mathrm{P}\)（\(\mathrm{K}\) 的一个为完满域的扩张），使得 \(\mathrm{A}_{(\mathrm{P})}\) -模 \(\mathrm{M}_{(\mathrm{P})}\) 是半单的。
\item
  A-模 M 是半单的，且其交换子的中心是域 K 上的一个平展代数。
\end{enumerate}

显然 (i) 蕴含 (ii)。假设 (ii) 成立；则由 VIII, 第 222 页的推论 1, a)，\(A_{(P)}\) -模 \(M_{(P)}\) 是绝对半单的。由命题 1, b)，M 于是是绝对半单的。因此 (ii) 蕴含 (i)。

假设 M 是半单的。设 Z 是 M 的交换子的中心；它是域 K 上的一个有限次数交换代数。若 L 是 K 的一个扩张，则由 VIII, 第 222 页的命题 8, b) 得出，\(A_{(L)}\) -模 \(M_{(L)}\) 是半单的当且仅当环 \(Z_{(L)}\) 是约化的。因此 (i) 与 (iii) 的等价性由 V, §6, No.~7, 第 34 页的定理 4 得出。

\subsection{可分闭域上的代数}

引理 1. —— 设 D 是一个体，Z 是它的中心。用 p 表示 Z 的特征指数。假设对 D 的每个元素 a，存在整数 \(m \geqslant 0\) 使得 \(a^{p^{m}}\) 属于 Z。则体 D 是交换的。

若 p = 1，则 D = Z。因此我们假设 p \textgreater{} 1。

我们用反证法来证明，假设 D 不是交换的。设 \(a\) 是 \(\mathrm{D} - \mathrm{Z}\) 的一个元素，\(q\) 是 \(p\) 的一个幂，使得 \(a^q\) 属于 Z。用 I 表示 D 上的恒等映射，把内自同构 \(x \mapsto axa^{-1}\)（从 D 到 D、与 \(a\) 相伴）记作 \(\sigma\)；我们有 \(\sigma^q = \mathrm{I}\)，因为 \(a^q\) 属于 Z。我们有 \(\sigma - \mathrm{I} \neq 0\)，因为 \(a\) 不属于 Z；我们还有 \((\sigma - \mathrm{I})^q = \sigma^q - \mathrm{I} = 0\)，因为 Z 的特征为 \(p\)。设 \(f\) 是使 \((\sigma - \mathrm{I})^f \neq 0\) 的最大自然数；我们有 \(f \geqslant 1\)。选取 D 的一个元素 \(c\)，使得 \((\sigma - \mathrm{I})^f(c) \neq 0\)，并令

\[
x = (\sigma - \mathrm{I}) ^ {f - 1} (c), \qquad y = (\sigma - \mathrm{I}) (x) = (\sigma - \mathrm{I}) ^ {f} (c).
\]

按构造，我们有 \(y \neq 0\) 且 \(\sigma(y) = y\)；若令 \(z = y^{-1}x\)，则由此得出

\[
\sigma (z) = \sigma (y) ^ {- 1} \sigma (x) = y ^ {- 1} (y + x) = 1 + z
\]

因此 \(\sigma(z^{p^j}) = 1 + z^{p^j}\) 对每个自然数 \(j\) 成立。选取整数 \(m \geqslant 0\) 使得 \(z^{p^m}\) 属于 D 的中心 Z；我们有

\[
z ^ {p ^ {m}} = a z ^ {p ^ {m}} a ^ {- 1} = \sigma (z ^ {p ^ {m}}) = 1 + z ^ {p ^ {m}}.
\]

这一矛盾便导出引理 1。

命题 3. —— 设 K 是一个可分闭域（V, §7, No.~8, 第 45 页，定义 4），D 是 K 上的一个为体的有限次数代数。则 D 是交换的。

用 p 表示 K 的特征指数。设 a 是 D 的一个元素。环 \(K[a]\) 是 K 的一个代数扩张（V, §3, No.~1, 第 17 页，推论 1）。由于域 K 是可分闭的，由 V, §7, No.~7, 第 44 页的命题 13 得出，代数 \(K[a]\) 是 K 的一个 p-根式扩张。于是存在整数 \(m \geqslant 0\) 使得 \(a^{p^{m}}\) 属于 K。由引理 1，体 D 因而是交换的。

推论. —— 设 K 是一个可分闭域，A 是 K 上的一个半单的有限次数代数。则存在整数 \(r \geqslant 0\)、严格正整数 \(n_{1}, \ldots, n_{r}\) 以及 K 的有限次数扩张 \(K_{1}, \ldots, K_{r}\)，使得 A 同构于代数 \(\prod_{i=1}^{r} \mathbf{M}_{n_{i}}(\mathbf{K}_{i})\)。

由半单代数的结构定理（VIII, 第 135 页，定理 1），A 同构于一个代数 \(\prod_{i=1}^{r} \mathbf{M}_{n_i}(\mathrm{D}_i)\)，其中 \(r\) 是整数 \(\geqslant 0\)，\(n_1, \ldots, n_r\) 是严格正整数，而 \(\mathrm{D}_1, \ldots, \mathrm{D}_r\) 是 K 上的为体的有限次数代数。由于域 K 是可分闭的，由命题 3，每个体 \(\mathrm{D}_i\) 都是交换的；推论得证。

\subsection{绝对半单代数}

定义 2. —— 设 K 是一个交换域。我们称 K-代数 A 是绝对半单的，如果对 K 的每个扩张 L，环 \(\mathrm{A}_{(\mathrm{L})}\) 都是半单的。

绝对半单代数是半单的。K-代数 A 绝对半单当且仅当 A-模 \(A_{s}\) 绝对半单。于是，由 VIII, 第 229 页的命题 1，我们得到如下结果：若

L 是 K 的一个扩张，则 L-代数 \(A_{(L)}\) 绝对半单当且仅当 K-代数 A 绝对半单。

定理 1. —— 设 K 是一个交换域，A 是一个 K-代数。下列性质等价：

\begin{enumerate}
\def\labelenumi{(\roman{enumi})}
\item
  K-代数 A 是绝对半单的。
\item
  代数 \(\mathrm{A}\) 在 \(\mathrm{K}\) 上的次数有限，且存在 \(\mathrm{P}\)（\(\mathrm{K}\) 的一个为完满域的扩张），使得 \(\mathrm{P}\) -代数 \(\mathrm{A}_{(\mathrm{P})}\) 是半单的。
\item
  K-代数 A 是半单的且在 K 上的次数有限，并且其中心是一个平展 K-代数。
\item
  存在一个有限族 \((n_{i},\mathrm{D}_{i})_{i\in\mathrm{I}}\)，其中 \(n_{i}\) 是严格正整数，而 \(D_{i}\) 是 K 上的一个为体的有限次数代数，使得 \(Z_{i}\)（即 \(D_{i}\) 的中心）对每个 \(i\in I\) 都是 K 的一个可分扩张，且 A 同构于矩阵环 \(\mathbf{M}_{n_{i}}(\mathrm{D}_{i})\) 之积。
\item
  存在 \(\mathrm{L}\)（\(\mathrm{K}\) 的一个扩张）与一个有限整数族 \((n_i)_{i\in \mathrm{I}}\)，使得 L-代数 \(\mathrm{A}_{(\mathrm{L})}\) 同构于代数 \(\prod_{i\in \mathrm{I}}\mathbf{M}_{n_i}(\mathrm{L})\)。
\item
  存在 K 的一个有限次数伽罗瓦扩张 L 与一个有限整数族 \((n_{i})\)，使得 \(\mathrm{A}_{(\mathrm{L})}\) 同构于矩阵代数 \(\mathbf{M}_{n_{i}}(\mathrm{L})\) 之积。
\end{enumerate}

我们先证明诸蕴含关系 (v) \(\Rightarrow\) (iv) \(\Rightarrow\) (iii) \(\Rightarrow\) (ii) \(\Rightarrow\) (i)。用 Z 表示 A 的中心。

若性质 (v) 成立，则 L-代数 \(\mathrm{A}_{(\mathrm{L})}\) 是半单的且在 L 上的次数有限，其中心（同构于 \(\mathrm{Z}_{(\mathrm{L})}\)，见 III, 第 41 页的推论）同构于 \(L^{r}\)，其中 \(r \geqslant 0\) 是某个整数。由 VIII, 第 222 页的推论 2, a)，代数 A 是半单的且在 K 上的次数有限。因此它同构于环的一个有限乘积 \(\prod_{i \in I} M_{n_{i}}(D_{i})\)，其中 \(n_{i} \geqslant 1\) 对每个 \(i \in I\) 成立，而体 \(D_{i}\) 是 K 上的一个有限次数代数。\(Z_{i}\)（即 \(D_{i}\) 的中心）是一个交换域且是 K 的一个扩张，而 Z 同构于 \(\prod_{i \in I} Z_{i}\)。于是，\(\mathrm{Z}_{(\mathrm{L})}\) 一方面同构于 \(\prod_{i \in I} (\mathrm{Z}_{i})_{(\mathrm{L})}\)，另一方面同构于 \(L^{r}\)。换言之，代数 \(\prod_{i \in I} Z_{i}\) 在域 K 上是平展的（V, §6, No.~3, 第 28 页，定义 1），且诸扩张 \(Z_{i}\) 在 K 上都是可分的（V, §6, No.~4, 第 31 页，推论，及 V, §7, No.~1, 第 42 页）。故 (v) 蕴含 (iv)。

若性质 (iv) 成立，则显然 A 是一个半单代数且在 K 上的次数有限。其中心 Z 同构于 K 的有限次数可分扩张的乘积 \(\prod_{i\in I}Z_{i}\)；因而它是一个平展代数（同前所引）。故 (iv) 蕴含 (iii)。

蕴含关系 (iii) \(\Rightarrow\) (ii) \(\Rightarrow\) (i) 由应用到 A-模 \(\mathrm{A}_s\) 的 VIII, 第 230 页的命题 2 得出。

引理 2. —— 设 L 是一个代数闭域，D 是一个在其中心中包含 L 的体。若 D 异于 L，则存在 L 的一个扩张 \(L'\)，使得环 \(D \otimes_{L} L'\) 不是左阿廷的。

设 \(x\) 是 \(\mathrm{D} - \mathrm{L}\) 的一个元素；由于 \(\mathrm{L}\) 是代数闭的，扩张 \(\mathrm{L}' = \mathrm{L}(x)\)（\(\mathrm{L}\) 的扩张）不是代数的，且 \(x\) 在 \(\mathrm{L}\) 上是超越的。于是环 \(\mathrm{B} = \mathrm{L}' \otimes_{\mathrm{L}} \mathrm{L}'\) 是一个整环，这由 \(\mathrm{V}\), §17, No.~4, 第 141 页的命题 5 得出。

B 的元素 \(y = x \otimes 1 - 1 \otimes x\) 不是零，但若 \(\varphi\) 是从 B 到 \(L'\) 的把 \(\xi \otimes \eta\) 映为 \(\xi \eta\) 的同态，则我们有 \(\varphi(y) = 0\)，故 y 在 B 中不可逆。我们把环 \(C = D \otimes_{L} L'\) 看作它的子环 B 上的右模；它是一个自由模，因为 D 是它的子域 \(L'\) 上的一个右向量空间。由于 y 是整环 B 的非零且不可逆的元素，在 C 中用 y 作右乘是一个单但非双射的映射 \(R_y\)。而 \(R_y\) 是左 C-模 \(C_s\) 的一个自同态；因此（VIII, 第 28 页，推论 1），环 C 不是左阿廷的。

我们现在来证明 (i) 蕴含 (v)。这是下述引理的一个推论。

引理 3. —— 设 A 是一个绝对半单 K-代数，L 是 K 的一个代数闭扩张。则代数 \(\mathrm{A}_{(\mathrm{L})}\) 同构于 L 上有限多个矩阵代数的一个乘积。

L-代数 \(\mathrm{A}_{(\mathrm{L})}\) 是半单的；因此它同构于有限多个形如 \(\mathbf{M}_{n_{i}}(\mathrm{D}_{i})\) 的代数的乘积，其中 \(D_{i}\) 是一个在其中心中包含 L 的体，而 \(n_{i}\) 是整数 \(\geqslant 1\)（VIII, 第 137 页，注 1）。

设 \(L'\) 是 L 的一个扩张。由于 K-代数 A 是绝对半单的，环 \(A_{(L')}\) 是半单的，因而是左阿廷的。现在，环 \(A_{(L')}\) 同构于 \(L' \otimes_L A_{(L)}\)，从而同构于 \(\prod_{i \in I} M_{n_i}(L' \otimes_L D_i)\)；由 VIII, 第 7 页的命题 5，诸环 \(M_{n_i}(L' \otimes_L D_i)\) 因而都是左阿廷的。

设 \(n \geqslant 1\) 是一个整数，B 是一个环。设 \((\mathfrak{b}_r)_{r \geqslant 0}\) 是 B 的左理想的一个递降序列；用 \(\mathfrak{c}_r\) 表示 \(n\) 阶且元素在 \(\mathfrak{b}_r\) 中的方阵所成之集。则 \((\mathfrak{c}_r)_{r \geqslant 0}\) 是 \(\mathbf{M}_n(\mathrm{B})\) 的左理想的一个递降序列。特别地，若环 \(\mathbf{M}_n(\mathrm{B})\) 是左阿廷的，则 B 也是。

由上所述，对每个 \(i \in I\) 与 L 的每个扩张 \(L'\)，环 \(D_{i} \otimes_{L} L'\) 都是左阿廷的。由引理 2，我们有 \(D_{i} = L\)（对每个 \(i \in I\)），由此即得引理 3。

我们将用下面的引理来证明蕴含关系 (i)⇒(vi)。

引理 4. —— 设 A 与 B 是域 K 上的具有有限生成系的代数，\(K'\) 是 K 的一个扩张。若 \(K'\) -代数 \(\mathrm{A}_{(\mathrm{K}')}\) 与 \(\mathrm{B}_{(\mathrm{K}')}\) 同构，则存在 \(K'\) 的一个在 K 上有限生成的子扩张 L，使得 L-代数 \(\mathrm{A}_{(\mathrm{L})}\) 与 \(\mathrm{B}_{(\mathrm{L})}\) 同构。

设 \((e_{i})_{i\in\mathrm{I}}\) 与 \((f_{j})_{j\in\mathrm{J}}\) 分别是代数 A 与 B 的有限生成系。设 u 是从 \(\mathrm{A}_{(\mathrm{K}^{\prime})}\) 到 \(\mathrm{B}_{(\mathrm{K}^{\prime})}\) 的一个同构，v 是其逆同构；存在 \(K^{\prime}\) 的一个在 K 上有限生成的子扩张 L，使得 \(u(1\otimes e_{i})\in\mathrm{B}_{(\mathrm{L})}\)（对每个 \(i\in\mathrm{I}\)）且 \(v(1\otimes f_{j})\in\mathrm{A}_{(\mathrm{L})}\)（对每个 \(j\in\mathrm{J}\)）。于是，u 把 \(\mathrm{A}_{(\mathrm{L})}\) 映到 \(\mathrm{B}_{(\mathrm{L})}\) 中，而 v 把 \(\mathrm{B}_{(\mathrm{L})}\) 映到 \(\mathrm{A}_{(\mathrm{L})}\) 中。诱导映射 \(u^{\prime}:\mathrm{A}_{(\mathrm{L})}\to\mathrm{B}_{(\mathrm{L})}\) 与 \(v^{\prime}:\mathrm{B}_{(\mathrm{L})}\to\mathrm{A}_{(\mathrm{L})}\) 是环同态；它们是互逆的双射。

我们来完成蕴含关系 \((\mathrm{i})\Rightarrow(\mathrm{vi})\) 的证明。设 \(K'\) 是 K 的一个可分闭包（V, §7, No.~8, 第 45 页）。则 \(A_{(K')}\) 是 \(K'\) 上的一个绝对半单代数。由蕴含关系 \((\mathrm{i})\Rightarrow(\mathrm{iv})\)，\(K'\) -代数 \(A_{(K')}\) 同构于一个乘积 \(\mathbf{M}_{n_{1}}(\mathrm{D}_{1})\times\cdots\times\mathbf{M}_{n_{r}}(\mathrm{D}_{r})\)，其中 \(D_{i}\) 是 \(K'\) 上的一个为体的有限次数代数，其中心 \(Z_{i}\) 是 \(K'\) 的一个可分扩张。由于 \(K'\) 是可分闭的，我们有 \(Z_{i}=K'\)。由 VIII, 第 231 页的命题 3，体 \(D_{i}\) 是交换的。因此我们有 \(D_{i}=K'\)。用 B 表示 K-代数 \(\mathbf{M}_{n_{1}}(\mathrm{K})\times\cdots\times\mathbf{M}_{n_{r}}(\mathrm{K})\)。由上所述，\(K'\) -代数 \(A_{(K')}\) 与 \(B_{(K')}\) 同构。\(K'\) 的每个在 K 上有限生成的子扩张都是可分的且在 K 上的次数有限，因此，每一个这样的子扩张都包含于某个子扩张 L 之中，其中 L 是 \(K'\) 的在 K 上伽罗瓦且次数有限的子扩张（V, §10, No.~1, 第 57 页，命题 2）。于是该蕴含关系由引理 4 得出。

蕴含关系 (vi)⇒(v) 是直接的。

推论 1. —— 设 K 是一个交换域，\(A_{1}\) 与 \(A_{2}\) 是两个 K-代数。假设 \(A_{1}\) 是绝对半单的。

\begin{enumerate}
\def\labelenumi{\alph{enumi})}
\item
  若 \(\mathrm{A}_2\) 是半单的，则 \(\mathrm{A}_1\otimes_{\mathrm{K}}\mathrm{A}_2\) 亦然
\item
  若 \(\mathrm{A}_2\) 是绝对半单的，则 \(\mathrm{A}_1\otimes_{\mathrm{K}}\mathrm{A}_2\) 亦然。
\end{enumerate}

把 \(A_{1}\) 的中心记作 \(Z_{1}\)，把 \(A_{2}\) 的中心记作 \(Z_{2}\)。由 III, §4, No.~4, 第 468 页的推论，\(A_{1} \otimes_{K} A_{2}\) 的中心 Z 等于 \(Z_{1} \otimes_{K} Z_{2}\)。假设 \(A_{2}\) 是半单的；则 \(Z_{2}\) 是一个约化代数（VIII, 第 137 页，命题 2 与 3）。由定理 1，\(Z_{1}\) 是平展的，因而是可分的

K-代数。由 V, §15, No.~2, 第 120 页的命题 5，环 \(Z = Z_{1} \otimes_{K} Z_{2}\) 是约化的；由于 \(A_{1}\) 在 K 上的次数有限（定理 1），由 VIII, 第 221 页的命题 7 得出，环 \(A_{1} \otimes_{K} A_{2}\) 是半单的。

现在假设 \(A_{2}\) 是绝对半单的。设 L 是 K 的一个扩张。则代数 \(\mathrm{A}_{1(\mathrm{L})}\) 是绝对半单的，而代数 \(\mathrm{A}_{2(\mathrm{L})}\) 是半单的。因此由 a)，代数 \(\mathrm{A}_{1} \otimes_{\mathrm{K}} \mathrm{A}_{2(\mathrm{L})}\) 是半单的。

推论 2. —— 设 K 是一个可分闭域，A 是一个绝对半单 K-代数。则存在整数 \(r \geqslant 0\) 与严格正整数 \(n_{1}, \ldots, n_{r}\)，使得代数 A 同构于代数 \(\prod_{i=1}^{r} \mathbf{M}_{n_{i}}(\mathrm{K})\)。

由定理 1，A 同构于形如 \(\prod_{i=1}^{r} \mathbf{M}_{n_i}(\mathrm{D}_i)\) 的一个代数，其中 \(r \geqslant 0\) 是整数，\(n_1, \ldots, n_r\) 是整数，而 \(\mathrm{D}_1, \ldots, \mathrm{D}_r\) 是 K 上的为体的有限次数代数，它们的中心是 K 的可分扩张，因而等于 K。由 VIII, 第 231 页的命题 3，我们有 \(\mathrm{D}_i = \mathrm{K}\)（对 \(i \in [1, r]\)）。

例. —— 一个交换 K-代数是绝对半单的当且仅当它是平展的：这由定义（V, §6, No.~3, 第 28 页，定义 1）与定理 1 的性质 (i) 与 (v) 的等价性得出。

\subsection{绝对半单模的刻画}

命题 4. —— 设 K 是一个交换域，A 是一个 K-代数。

\begin{enumerate}
\def\labelenumi{\alph{enumi})}
\item
  设 M 是一个半单 A-模。A-模 M 绝对半单当且仅当属于 M 的支集的每个单模都是绝对半单的。
\item
  设 S 是一个单 A-模，D 是它的交换子。下列性质等价：
\end{enumerate}

\begin{enumerate}
\def\labelenumi{(\roman{enumi})}
\item
  A-模 S 是绝对半单的。
\item
  K-代数 \(\mathbf{D}\) 是绝对半单的。
\item
  K-代数 D 是一个体，在 K 上的次数有限，且其中心是 K 的一个可分扩张。
\end{enumerate}

断言 a) 由 VIII, 第 229 页的命题 1, a) 得出。设 S 与 D 如 b) 中所述，并设 L 是 K 的一个扩张。\(A_{(L)}\) -模 \(S_{(L)}\) 是半单的当且仅当环 \(D_{(L)}\) 是半单的（VIII, 第 222 页，命题 8, c)）。这就证明了 (i) 与 (ii) 的等价性，而 (ii) 与 (iii) 的等价性由定理 1 得出，因为 D 是一个体。

推论. —— 设 K 是一个交换域，\(A_{1}\) 与 \(A_{2}\) 是两个 K-代数，\(M_{1}\) 是一个绝对半单 \(A_{1}\) -模，而 \(M_{2}\) 是一个半单 \(A_{2}\) -模。则 \(M_{1} \otimes_{K} M_{2}\) 是环 \(A_{1} \otimes_{K} A_{2}\) 上的一个半单模。

模 \(M_{1}\) 是绝对半单的单 \(A_{1}\) -模的直和（命题 4）。因此只需在模 \(M_{1}\) 与 \(M_{2}\) 都是单的情形证明断言。用 \(D_{1}\) 与 \(D_{2}\) 表示它们的交换子。K-代数 \(D_{1}\) 是绝对半单的（同前所引）；由 VIII, 第 234 页的推论 1，K-代数 \(D_{1} \otimes_{K} D_{2}\) 是半单的。于是由 VIII, 第 215 页的推论 1 得出，\((\mathrm{A}_{1} \otimes_{\mathrm{K}} \mathrm{A}_{2})\) -模 \(M_{1} \otimes_{K} M_{2}\) 是半单的。

\subsection{半单代数上的导子}

在本小节及随后各小节中，K 是一个交换环，A 是一个 K-代数，B 是 K-代数 \(A \otimes_{K} A^{o}\)，而 \(\varepsilon\) 是由 \(\varepsilon(x \otimes y) = xy\)（对 A 中的 x, y）定义的从 B 到 A 的 K-线性映射。

回顾（III, §4, No.~3, 第 467 页）每个 (A, A)-双模 P 都可看作一个左 B-模，其外部作用律由 \((a \otimes a')p = apa'\)（对 A 中的 \(a, a'\) 与 P 中的 p）给出。反之，每个 B-模都可看作一个 (A, A)-双模。我们赋予 A 它的典范 (A, A)-双模结构及相应的 B-模结构；我们赋予 B 对应于 B-模 \(B_s\) 的 (A, A)-双模结构。于是我们有

\[
a (x \otimes y) a ^ {\prime} = (a \otimes a ^ {\prime}) (x \otimes y) = a x \otimes y a ^ {\prime}
\]

对 A 中的 \(a, a', x, y\) 成立，其中乘积 \(ya'\) 是在代数 A 中计算的。

K-线性映射 \(\varepsilon\) 是 (A, A)-双模的一个同态。

命题 5. —— 下列性质等价：

\begin{enumerate}
\def\labelenumi{(\roman{enumi})}
\item
  B-模 A 是投射的。
\item
  存在 (A,A)-双模 B 的一个元素 \(e\)，满足以下两个条件：\(\varepsilon(e) = 1\) 且 \(ae = ea\)（对每个 \(a \in \mathrm{A}\)）。
\end{enumerate}

映射 \(\varepsilon: \mathrm{B} \to \mathrm{A}\) 是满射的，因为 \(\varepsilon(a \otimes 1) = a\) 对每个 \(a \in \mathrm{A}\) 成立；它是 (A, A)-模的一个同态，因而是 B-线性的。若 B-模 A 是投射的，则存在一个截面 \(s\)（\(\varepsilon\) 的；见 II, §2, No.~2, 第 231 页，命题 4）；它是从 A 到 B 的一个 (A, A)-双模同态。若令 \(e = s(1)\)，则我们有 \(\varepsilon(e) = \varepsilon(s(1)) = 1\) 与 \(ae = s(a) = ea\)（对每个 \(a \in \mathrm{A}\)）。故 (i) 蕴含 (ii)。

反之，设 e 是 B 的一个满足 (ii) 中条件的元素。用下述公式定义从 A 到 B 的一个映射 s：

\[
s (a) = a e = e a.\tag{1}
\]

它是 (A, A)-模的一个同态，且我们有 \(\varepsilon \circ s = 1_{\mathrm{A}}\)；换言之，\(s\) 是满射映射 \(\varepsilon\) 的一个 B-线性截面。因此，B-模 A 同构于子模 \(s(\mathrm{A})\)（\(\mathrm{B}_s\) 的直因子子模；见 II, §1, No.~9, 第 211 页，命题 15），从而是投射的（II, §2, No.~2, 第 231 页，命题 4）。这就证明了 (ii) 蕴含 (i)。

注. —— 1) 设 \(e = \sum_{i=1}^{r} a_i \otimes a_i'\) 是 B 的一个元素。命题 5 的 (ii) 中的条件翻译为公式

\[
\sum_ {i = 1} ^ {r} a _ {i} a _ {i} ^ {\prime} = 1,\tag{2}
\]

\[
\sum_ {i = 1} ^ {r} a a _ {i} \otimes a _ {i} ^ {\prime} = \sum_ {i = 1} ^ {r} a _ {i} \otimes a _ {i} ^ {\prime} a \quad \text { for   every } a \in A.\tag{3}
\]

当它们被满足时，e 是 B 中的一个幂等元。事实上，我们这时有关系式

\[
e ^ {2} = \sum_ {i = 1} ^ {r} a _ {i} e a _ {i} ^ {\prime} = \sum_ {i = 1} ^ {r} e a _ {i} a _ {i} ^ {\prime} = e.
\]

\begin{enumerate}
\def\labelenumi{\arabic{enumi})}
\setcounter{enumi}{1}
\tightlist
\item
  设 K 是一个交换域，A 是一个 K-代数，M 是一个 A-模。加群 \(\operatorname{End}_{\mathrm{K}}(\mathrm{M})\) 被赋予一个由下式定义的 (A, A)-双模结构
\end{enumerate}

\[
a u a ^ {\prime} (x) = a u \left(a ^ {\prime} x\right)
\]

对每个 \(a, a' \in A\)、每个 \(u \in \operatorname{End}_{\mathrm{K}}(\mathrm{M})\) 与每个 \(x \in M\) 成立。我们赋予它相应的 B-模结构。设 \(e = \sum_{i=1}^{r} a_i \otimes a'_i\) 是 B 的一个满足命题 5 中 (ii) 的条件的元素，因而也满足关系式 (2) 与 (3)。若 \(p \in \operatorname{End}_{\mathrm{K}}(\mathrm{M})\) 是一个投影，其像 N 是 M 的一个 A-子模，则 ep 是具有同一像的一个 A-线性投影。

事实上，\(ep\) 的像包含在 N 中。若 \(x\) 属于 N，则 \(a_i' x\) 亦然，于是我们有 \(p(a_i' x) = a_i' x\) 且

\[
e p (x) = \sum_ {i = 1} ^ {r} a _ {i} a _ {i} ^ {\prime} x = x
\]

这由公式 (2) 得出。由公式 (3) 我们推出，\(aep(x) = ep(ax)\) 对每个 \(a \in \mathrm{A}\) 与每个 \(x \in \mathrm{M}\) 成立，这就证明了 \(ep\) 是 A-线性的。

定理 2. —— 设 K 是一个交换域，A 是一个 K-代数。下列性质等价：

\begin{enumerate}
\def\labelenumi{(\roman{enumi})}
\item
  K-代数 A 是绝对半单的。
\item
  K-代数 \(\mathrm{B} = \mathrm{A}\otimes_{\mathrm{K}}\mathrm{A}^{\circ}\) 是半单的。
\item
  B-模 A 是投射的。
\item
  存在 (A,A)-双模 B 的一个元素 \(e\)，满足 \(\varepsilon(e) = 1\) 且 \(ae = ea\)（对每个 \(a \in \mathrm{A}\)）。
\end{enumerate}

假设代数 A 是绝对半单的，因而是半单的。则代数 \(A^{o}\) 是半单的（VIII, 第 137 页，命题 2），且由 VIII, 第 234 页的推论 1 得出，\(B = A \otimes_{K} A^{o}\) 是一个半单 K-代数。因此，(i) 蕴含 (ii)。

由于半单环上的每个模都是投射的（VIII, 第 138 页，命题 4），(ii) 蕴含 (iii)。

(iii) 与 (iv) 的等价性由命题 5 得出。为完成证明，我们来证明 (iv) 蕴含 (i)。设 \(e = \sum_{i=1}^{r} a_i \otimes a_i'\) 是 B 的一个满足命题 5 中 (ii) 的条件的元素。设 L 是域 K 的一个扩张；我们必须证明环 \(\mathrm{A}_{(\mathrm{L})}\) 是半单的，或者等价地，每个 \(\mathrm{A}_{(\mathrm{L})}\) -模都是半单的（VIII, 第 138 页，命题 4）。设 M 是一个 \(\mathrm{A}_{(\mathrm{L})}\) -模，N 是 M 的一个子模；我们把 M 看作左 (A,L)-双模，把 N 看作一个子双模（III, §4, No.~3, 第 466 页）。由于 L 是一个域，存在 M 中以 N 为像的 L-线性投影 \(u\)。由于位似 \(a_{\mathrm{M}}\)（与 A 的元素 \(a\) 相伴的）都是 L-线性的，存在从 \(\mathrm{A} \otimes_{\mathrm{K}} \mathrm{A}^{\circ}\) 到 \(\operatorname{End}_{\mathrm{L}}(\mathrm{M})\) 的唯一的加群同态，它把元素 \(a \otimes a'\) 映为 L-线性映射 \(x \mapsto au(a'x)\)。我们把此同态下 \(e\) 的像记作 \(v\)；由注 2 得出，\(v\) 是一个以 N 为像的 \(\mathrm{A}_{(\mathrm{L})}\) -线性投影。\(v\) 的核是 M 的一个 \(\mathrm{A}_{(\mathrm{L})}\) -子模，它与 N 互补。由 VIII, 第 56 页的推论 2，\(\mathrm{A}_{(\mathrm{L})}\) -模 M 是半单的。

注 3. —— 我们知道（VIII, 第 234 页，推论 1），交换域上的两个绝对半单代数的张量积是绝对半单的。因此，若代数 A 是绝对半单的，则代数 \(B = A \otimes_{K} A^{\circ}\) 亦然。

\subsection{代数的上同调}

在本小节中，K 是一个交换环，A 是一个 K-代数，B 是 K-代数 \(A \otimes_{K} A^{\circ}\)，而 \(\varepsilon\) 是由 \(\varepsilon(x \otimes y) = xy\)（对 A 中的 x, y）定义的从 B 到 A 的 K-线性映射。对 \(n \in N\)，我们把 K-模 A 的 \(n + 2\) 个拷贝在 K 上的张量积记作 \(B_{n}\)。我们把它看作一个 (A, A)-双模（也看作一个 B-模）；我们赋予它由张量积第一个因子的左 A-模结构导出的左 A-模结构，以及由最后一个因子的右 A-模结构导出的右 A-模结构。特别地，\(B_{0}\) 就是 (A, A)-双模 B。

对每个整数 \(n \geqslant 1\)，我们用下述公式定义双模同态 \(d_{n}^{i}\)（对 \(i \in [0, n]\)）与 \(d_{n}\)（从 \(B_{n}\) 到 \(B_{n-1}\)）

\[
d _ {n} ^ {i} (x _ {0} \otimes \dots \otimes x _ {n + 1}) = x _ {0} \otimes \dots \otimes x _ {i} x _ {i + 1} \otimes \dots \otimes x _ {n + 1}\tag{4}
\]

（对 \(i\in [0,n]\)）以及

\[
d _ {n} = \sum_ {i = 0} ^ {n} (- 1) ^ {i} d _ {n} ^ {i}.\tag{5}
\]

我们把映射 \(\varepsilon : \mathrm{B}_0 \to \mathrm{A}\) 记作 \(d_0 = d_0^0\)。

设 \(n\) 是一个整数 \(\geqslant 1\)。对 \(0 \leqslant i < j \leqslant n\)，我们有

\[
d _ {n - 1} ^ {i} \circ d _ {n} ^ {j} = d _ {n - 1} ^ {j - 1} \circ d _ {n} ^ {i},\tag{6}
\]

由此我们推出

\[
\begin{array}{c} d _ {n - 1} \circ d _ {n} = \sum_ {0 \leqslant i <   j \leqslant n} (- 1) ^ {i + j} d _ {n - 1} ^ {i} \circ d _ {n} ^ {j} + \sum_ {0 \leqslant j \leqslant i \leqslant n - 1} (- 1) ^ {i + j} d _ {n - 1} ^ {i} \circ d _ {n} ^ {j} \\ = \sum_ {0 \leqslant i <   j \leqslant n} (- 1) ^ {i + j} d _ {n - 1} ^ {j - 1} \circ d _ {n} ^ {i} + \sum_ {0 \leqslant j \leqslant i \leqslant n - 1} (- 1) ^ {i + j} d _ {n - 1} ^ {i} \circ d _ {n} ^ {j} \end{array}
\]

从而

\[
d _ {n - 1} \circ d _ {n} = 0.\tag{7}
\]

设 P 是一个 (A,A)-双模。对任意整数 \(n \geqslant 0\)，我们把从 \(A^{n}\) 到 P 的 K-多重线性映射所成的 K-模记作 \(C^{n}(A,P)\)。映射 \(\alpha^{n}: C^{n}(A,P) \to \operatorname{Hom}_{\mathrm{B}}(\mathrm{B}_{n},\mathrm{P})\) 把 \(f \in \mathrm{C}^{n}(A,P)\) 映为同态 \(\alpha^{n}(f)\)，后者由下式刻画

\[
\alpha^ {n} (f) \left(x _ {0} \otimes \dots \otimes x _ {n + 1}\right) = x _ {0} f (x _ {1}, \ldots , x _ {n}) x _ {n + 1}\tag{8}
\]

是 K-模的一个同构。

我们用 \(\partial^{n}\)（对 \(n \geqslant 0\)）表示从 \(\mathrm{C}^{n}(\mathrm{A},\mathrm{P})\) 到 \(\mathrm{C}^{n+1}(\mathrm{A},\mathrm{P})\) 的使下图交换的唯一的 K-线性映射：

\[
\begin{array}{c} \mathrm{C} ^ {n} (\mathrm{A}, \mathrm{P}) \xrightarrow {\partial^ {n}} \mathrm{C} ^ {n + 1} (\mathrm{A}, \mathrm{P}) \\ \Biggl \downarrow_ {\alpha^ {n}} \\ \mathrm{Hom} _ {\mathrm{B}} (\mathrm{B} _ {n}, \mathrm{P}) \xrightarrow {\mathrm{Hom} (d _ {n + 1} , 1 _ {\mathrm{P}})} \mathrm{Hom} _ {\mathrm{B}} (\mathrm{B} _ {n + 1}, \mathrm{P}). \end{array}
\]

因此，按定义我们有

\[
(\alpha^ {n + 1} \circ \partial^ {n}) (f) = \alpha^ {n} (f) \circ d _ {n + 1}\tag{9}
\]

对每个 \(f \in \mathrm{C}^n(\mathrm{A},\mathrm{P})\)。换言之，我们有

\[
\partial^ {n} (f) \left(x _ {0}, \dots , x _ {n}\right) = \alpha^ {n} (f) \left(d _ {n + 1} \left(1 \otimes x _ {0} \otimes \dots \otimes x _ {n} \otimes 1\right)\right)
\]

对 A 中的 \(x_0, \ldots, x_n\) 以及 \(f\)（在 \(\mathrm{C}^n(\mathrm{A}, \mathrm{P})\) 中）成立，也就是说，

\[
\begin{array}{r l} & {\partial^ {n} (f) (x _ {0}, \ldots , x _ {n}) = x _ {0} f (x _ {1}, \ldots , x _ {n})} \\ & {\qquad + \sum_ {i = 0} ^ {n - 1} (- 1) ^ {i + 1} f (x _ {0}, \ldots , x _ {i - 1}, x _ {i} x _ {i + 1}, x _ {i + 2}, \ldots , x _ {n})} \\ & {\qquad \qquad \qquad \qquad \qquad + (- 1) ^ {n + 1} f (x _ {0}, \ldots , x _ {n - 1}) x _ {n}.} \end{array}\tag{10}
\]

由 (7) 与 (9)，我们有

\[
\partial^ {n + 1} \circ \partial^ {n} = 0 \quad \text {   for   every   } n \geqslant 0.\tag{11}
\]

我们把 K-模 \(\operatorname{Ker} \partial^0\) 记作 \(\mathrm{H}^0(\mathrm{A},\mathrm{P})\)，而对 \(n \geqslant 1\)，把 K-模 \(\operatorname{Ker} \partial^n/\operatorname{Im} \partial^{n-1}\) 记作 \(\mathrm{H}^n(\mathrm{A},\mathrm{P})\)。我们把 K-模 \(\mathrm{C}^0(\mathrm{A},\mathrm{P})\) 与 \(\mathrm{P}\) 等同起来，并且我们有 \(\mathrm{C}^1(\mathrm{A},\mathrm{P}) = \operatorname{Hom}_\mathrm{K}(\mathrm{A},\mathrm{P})\)。诸映射 \(\partial^n\)（对 \(n \leqslant 2\)）由下述公式给出

\[
\partial^ {0} (p) (a) = a p - p a \qquad \mathrm{for\ every} p \in \mathrm{P},\tag{12}
\]

\[
\partial^ {1} (f) \left(a, a ^ {\prime}\right) = a f \left(a ^ {\prime}\right) - f \left(a a ^ {\prime}\right) + f (a) a ^ {\prime} \quad \text { for } f \in \mathrm{C} ^ {1} (\mathrm{A}, \mathrm{P}),\tag{13}
\]

\[
\partial^ {2} (f) \left(a, a ^ {\prime}, a ^ {\prime \prime}\right) = a f \left(a ^ {\prime}, a ^ {\prime \prime}\right) - f \left(a a ^ {\prime}, a ^ {\prime \prime}\right) + f \left(a, a ^ {\prime} a ^ {\prime \prime}\right) - f \left(a, a ^ {\prime}\right) a ^ {\prime \prime}\tag{14}
\]

对 \(f \in \mathrm{C}^2(\mathrm{A},\mathrm{P})\) 成立。

于是 \(H^{0}(A,P)\) 是 P 的由满足对每个 \(a \in A\) 有 ap = pa 的元素 p 组成的 K-子模，而 \(H^{1}(A,P)\) 是从 A 到 P 的 K-导子所成的 K-模 \(\mathrm{Der}_{\mathrm{K}}(\mathrm{A},\mathrm{P})\)（III, §10, No.~2, 第 553 页）模去由形如 \(a \mapsto ap - pa\)（其中 \(p \in P\)）的导子组成的 K-子模所得的商（这些导子称为内导子）。

\subsection{绝对半单代数的上同调}

命题 6. —— 设 K 是一个交换环，A 是一个 K-代数。设 \(e = \sum_{i=1}^{r} a_i \otimes a_i'\) 是 \(B = A \otimes_K A^o\) 的一个满足 VIII, 第 236 页的命题 5 中 (ii) 的条件的元素。对任意整数 \(n \geqslant 1\) 与 \(C^n(A, P)\) 的任意元素 f，我们用 \(\gamma^n(f)\) 表示由下述公式定义的 \(C^{n-1}(A, P)\) 的元素

\[
\gamma^ {n} (f) \left(x _ {1}, \dots , x _ {n - 1}\right) = \sum_ {i = 1} ^ {r} a _ {i} f \left(a _ {i} ^ {\prime}, x _ {1} \dots , x _ {n - 1}\right).\tag{15}
\]

于是我们有

\[
\partial^ {n - 1} (\gamma^ {n} (f)) + \gamma^ {n + 1} (\partial^ {n} (f)) = f\tag{16}
\]

对每个整数 \(n \geqslant 1\) 与每个 \(f \in \mathrm{C}^n(\mathrm{A},\mathrm{P})\) 成立。

*注 1. —— 态射 \(\partial^n: \mathrm{C}^n(\mathrm{A},\mathrm{P}) \to \mathrm{C}^{n+1}(\mathrm{A},\mathrm{P})\) 定义了 K-模的一个复形 \((\mathrm{C}(\mathrm{A},\mathrm{P}),\partial)\)（X, §2, n° 1, 第 24 页）。于是映射 \(\gamma_n\) 定义了从该复形到它自身、把 0 连接到 \(\mathrm{Id}_{\mathrm{C}(\mathrm{A},\mathrm{P})}\) 的一个同伦（X, §2, n° 4, 第 32 页，定义 4）。 *

我们沿用 No.6 的记号。对每个整数 \(n \geqslant 0\)，我们用公式定义映射 \(h_n: \mathrm{B}_n \to \mathrm{B}_{n+1}\)

\[
h _ {n} (x) = d _ {n + 2} ^ {1} (e \otimes x) = \sum_ {i = 1} ^ {r} a _ {i} \otimes a _ {i} ^ {\prime} x.
\]

它是 \((A, A)\) -双模的一个同态（公式 (3)）。

引理 5. —— 我们有关系式

\[
d _ {n + 1} \circ h _ {n} + h _ {n - 1} \circ d _ {n} = 1 _ {\mathrm{B} _ {n}}\tag{17}
\]

对每个 \(n \geqslant 1\) 成立。

设 \(x\in \mathrm{B}_n\)；我们有

\[
\begin{array}{c} (d _ {n + 1} \circ h _ {n}) (x) = (d _ {n + 1} \circ d _ {n + 2} ^ {1}) (e \otimes x) \\ = (d _ {n + 1} ^ {0} \circ d _ {n + 2} ^ {1}) (e \otimes x) - \sum_ {i = 2} ^ {n + 2} (- 1) ^ {i} (d _ {n + 1} ^ {i - 1} \circ d _ {n + 2} ^ {1}) (e \otimes x), \end{array}
\]

于是，由公式 (6)，

\[
(d _ {n + 1} \circ h _ {n}) (x) = (d _ {n + 1} ^ {0} \circ d _ {n + 2} ^ {0}) (e \otimes x) - \sum_ {i = 2} ^ {n + 2} (- 1) ^ {i} (d _ {n + 1} ^ {1} \circ d _ {n + 2} ^ {i}) (e \otimes x).
\]

但我们有

\[
(d _ {n + 1} ^ {0} \circ d _ {n + 2} ^ {0}) (e \otimes x) = \varepsilon (e) x = x
\]

这由 VIII, 第 236 页的命题 5 的性质 (ii) 得出，而对 \(i \geqslant 2\)，

\[
d _ {n + 2} ^ {i} (e \otimes x) = e \otimes d _ {n} ^ {i - 2} (x),
\]

由此给出

\[
(d _ {n + 1} \circ h _ {n}) (x) = x - d _ {n + 1} ^ {1} (e \otimes d _ {n} (x)) = x - h _ {n - 1} \circ d _ {n} (x)
\]

从而得到公式 (17)。

利用引理 5，我们可以完成命题 6 的证明。设 n 是一个整数 \(\geqslant 1\)，f 是 \(C^{n}(A,P)\) 的一个元素。按构造，我们有

\[
\alpha^ {n - 1} (\gamma^ {n} (f)) = \alpha^ {n} (f) \circ h _ {n - 1}\tag{18}
\]

于是，由公式 (9) 与 (18)，

\[
\begin{array}{r l} \alpha^ {n} (\partial^ {n - 1} (\gamma^ {n} (f)) + \gamma^ {n + 1} (\partial^ {n} (f))) & = \alpha^ {n - 1} (\gamma^ {n} (f)) \circ d _ {n} + \alpha^ {n + 1} (\partial^ {n} (f)) \circ h _ {n} \\ & = \alpha^ {n} (f) \circ h _ {n - 1} \circ d _ {n} + \alpha^ {n} (f) \circ d _ {n + 1} \circ h _ {n} \\ & = \alpha^ {n} (f), \end{array}
\]

其中最后的等式由 (17) 得出。由于 \(\alpha^{n}\) 是双射，命题得证。

定理 3. —— 设 K 是一个交换环，A 是一个 K-代数，P 是一个 (A,A)-双模。假设 \((\mathrm{A}\otimes_{\mathrm{K}}\mathrm{A}^{\circ})\) -模 A 是投射的。则我们有 \(\mathrm{H}^{n}(\mathrm{A},\mathrm{P})=0\)（对每个整数 \(n\geqslant1\)）。

我们必须证明，对每个整数 \(n \geqslant 1\)，\(C^{n}(A, P)\) 的每个满足 \(\partial^{n}(f) = 0\) 的元素 f 都具有形式 \(\partial^{n-1}(g)\)（对 \(C^{n-1}(A, P)\) 的某个元素 g）。由命题 5，这是命题 6 的一个直接推论。

推论. —— 从 A 到 P 的每个 K-导子都是内导子。

这是等式 \(H^{1}(A,P)=0\) 的一个转述。

注. —— 2) 定理 3 的假设特别在 K 是一个域且 A 是一个绝对半单 K-代数时被满足（VIII, 第 238 页，定理 2）。

*3) 假设 K-模 A 是投射的。定理 3 也可按下述方式证明。复形 \((\bigoplus_{n\geqslant 0}\mathrm{B}_n,d)\) 与同态 \(\varepsilon :\mathrm{B}_0\to \mathrm{A}\) 定义了 B-模 A 的一个投射分解；因此，对每个 \(n\geqslant 0\)，K-模 \(\mathrm{H}^n (\mathrm{A},\mathrm{P})\) 同构于 \(\operatorname {Ext}_{\mathrm{B}}^{n}(\mathrm{A},\mathrm{P})\)（X, §6, n° 1, 第 100 页，定理 1）。若 B-模 A 是投射的，则 K-模 \(\operatorname {Ext}_{\mathrm{B}}^{n}(\mathrm{A},\mathrm{P})\) 对 \(n\geqslant 1\) 为零（X, §5, n° 3, 第 88 页，命题 5 的推论），由此推出 \(\mathrm{H}^n (\mathrm{A},\mathrm{P})\) 为零。反之，若对每个 (A, A)-双模 P 都有 \(\mathrm{H}^1 (\mathrm{A},\mathrm{P})\) 为零，则 B-模 A 是投射的（X, §5, n° 5, 第 93 页，命题 10）。*

\subsection{阿廷代数的分裂}

在本小节中，K 是一个交换环，A 是一个 K-代数。用 r 表示 A 的根。我们把商代数 \(A/r\) 记作 \(\overline{A}\)，把从 A 到 \(\overline{A}\) 的典范映射记作 \(\pi\)。我们关注的是 A 的满足 \(A = S \oplus r\) 的子代数 S。

我们用 \(\Sigma\) 表示 K-线性截面 \(s\)（\(\pi\) 的）中满足 \(s(\alpha\beta) = s(\alpha)s(\beta)\)（对 \(\alpha, \beta\) 属于 \(\overline{\mathrm{A}}\)）者所成之集。注意这样的截面必然满足 \(s(1) = 1\)（换言之，\(s\) 是一个环同态）：我们确有 \(s(1)^2 = s(1)\)，且 \(s(1)\) 是可逆的，因为它属于 \(1 + \mathfrak{r}\)（VIII, 第 156 页，定理 1）。若 \(s\) 是 \(\Sigma\) 的一个元素，则 \(s\) 的像 S 是 A 的一个子代数，且我们有 \(\mathrm{A} = \mathrm{S} \oplus \mathfrak{r}\)。反之，若 S 是 A 的一个满足 \(\mathrm{A} = \mathrm{S} \oplus \mathfrak{r}\) 的子代数，则 \(\pi\) 在 S 上的限制是双射的，且该逆双射定义了 \(\Sigma\) 的一个以 S 为像的元素。

由雅各布森定理（同前所引），\(1 + \mathfrak{r}\) 的每个元素在 A 中都是可逆的。我们称形如 \(a \mapsto xax^{-1}\)（其中 \(x \in 1 + \mathfrak{r}\)）的 A 的内自同构为特殊自同构。

命题 7. —— 假设 \((\overline{\mathrm{A}}\otimes_{\mathrm{K}}\overline{\mathrm{A}}^{\mathrm{o}})\) -模 \(\overline{\mathrm{A}}\) 是投射的。

\begin{enumerate}
\def\labelenumi{\alph{enumi})}
\item
  设 \(S_1\) 与 \(S_2\) 是 A 的满足 \(A = S_1 \oplus \mathfrak{r} = S_2 \oplus \mathfrak{r}\) 的子代数。存在 A 的一个把 \(S_1\) 变到 \(S_2\) 的特殊自同构。
\item
  假设 \(\pi\) 有一个 K-线性截面，且 A 的根 r 是幂零的。则存在 A 的一个满足 \(A = S \oplus r\) 的子代数 S。
\end{enumerate}

设 \(S_{1}\) 与 \(S_{2}\) 如 a) 中所述。设 \(s_{1}\) 与 \(s_{2}\) 是集合 \(\Sigma\) 中对应于子代数 \(S_{1}\) 与 \(S_{2}\) 的元素。设 \(\varepsilon\) 是从 \(A \otimes_{K} A\) 到 \(A\) 的由 \(\varepsilon(a \otimes b) = ab\) 给出的 K-线性映射。由 VIII, 第 236 页的命题 5 与 VIII, 第 237 页的注 1，存在元素 \(e = \sum_{i=1}^{r} \alpha_{i} \otimes \alpha_{i}'\)（属于 \(\overline{A} \otimes_{K} \overline{A}\)），满足 \(\sum_{i=1}^{r} \alpha_{i} \alpha_{i}' = 1\) 且 \(\sum_{i=1}^{r} \alpha \alpha_{i} \otimes \alpha_{i}' = \sum_{i=1}^{r} \alpha_{i} \otimes \alpha_{i}' \alpha\)（对每个 \(\alpha \in \overline{A}\)）。令 \(x = \sum_{i=1}^{r} s_{1}(\alpha_{i}) s_{2}(\alpha_{i}')\)。我们有 \(\pi(x) = \sum_{i=1}^{r} \alpha_{i} \alpha_{i}' = 1\)，因此 \(x \in 1 + r\)。设 \(\alpha\) 是 \(\overline{A}\) 的一个元素。我们有

\[
\begin{array}{l} s _ {1} (\alpha) x = \sum_ {i = 1} ^ {r} s _ {1} (\alpha \alpha_ {i}) s _ {2} (\alpha_ {i} ^ {\prime}) = (\varepsilon \circ (s _ {1} \otimes s _ {2})) \biggl (\sum_ {i = 1} ^ {r} \alpha \alpha_ {i} \otimes \alpha_ {i} ^ {\prime} \biggr) \\ = (\varepsilon \circ (s _ {1} \otimes s _ {2})) \biggl (\sum_ {i = 1} ^ {r} \alpha_ {i} \otimes \alpha_ {i} ^ {\prime} \alpha \biggr) = \sum_ {i = 1} ^ {r} s _ {1} (\alpha_ {i}) s _ {2} (\alpha_ {i} ^ {\prime} \alpha) = x s _ {2} (\alpha). \end{array}
\]

由此得到等式 \(x^{-1}S_{1}x = S_{2}\)，从而得到断言 a)。

在 b) 的假设下，再假设 \(r^{2}=0\)。在这种情形，(A, A)-双模 r 被 r 零化，因此我们把它看作一个 \((\overline{A},\overline{A})\) -双模。选取一个 K-线性截面 \(\sigma\)（\(\pi\) 的）。我们有

\[
\alpha x = \sigma (\alpha) x \quad \text { and } \quad x \alpha = x \sigma (\alpha)\tag{19}
\]

对 \(\alpha \in \overline{\mathrm{A}}\) 与 \(x\in \mathfrak{r}\) 成立。令

\[
\varphi (\alpha , \beta) = \sigma (\alpha \beta) - \sigma (\alpha) \sigma (\beta)\tag{20}
\]

对 \(\alpha, \beta \in \overline{\mathbf{A}}\) 成立。我们有关系式 \(\pi(\varphi(\alpha, \beta)) = \alpha\beta - \alpha\beta = 0\)（对 \(\alpha, \beta \in \overline{\mathbf{A}}\)）。因此，\(\varphi\) 定义了 \(\mathrm{C}^2(\overline{\mathbf{A}}, \mathfrak{r})\) 的一个元素。设 \(\alpha, \beta, \gamma\) 是 \(\overline{\mathbf{A}}\) 的元素；注意到 (19)，我们有

\[
\begin{array}{r l} \partial^ {2} \varphi (\alpha , \beta , \gamma) & = \alpha \varphi (\beta , \gamma) - \varphi (\alpha \beta , \gamma) + \varphi (\alpha , \beta \gamma) - \varphi (\alpha , \beta) \gamma \\ & = \sigma (\alpha) \varphi (\beta , \gamma) - \varphi (\alpha \beta , \gamma) + \varphi (\alpha , \beta \gamma) - \varphi (\alpha , \beta) \sigma (\gamma) \\ & = \sigma (\alpha) (\sigma (\beta \gamma) - \sigma (\beta) \sigma (\gamma)) - \sigma (\alpha \beta \gamma) + \sigma (\alpha \beta) \sigma (\gamma) + \sigma (\alpha \beta \gamma) \\ & - \sigma (\alpha) \sigma (\beta \gamma) - (\sigma (\alpha \beta) - \sigma (\alpha) \sigma (\beta)) \sigma (\gamma) \\ & = 0. \end{array}
\]

由 VIII, 第 242 页的定理 3，K-模 \(\mathrm{H}^{2}(\overline{\mathrm{A}},\mathfrak{r})\) 化为零。因此存在元素 \(\psi\)（\(\mathrm{C}^{1}(\overline{\mathrm{A}},\mathfrak{r})\) 的），使得 \(\partial^{1}\psi = \varphi\)，换句话说，使得我们有

\[
\varphi (\alpha , \beta) = \alpha \psi (\beta) - \psi (\alpha \beta) + \psi (\alpha) \beta \quad \text { for } \alpha , \beta \text { in } \overline {{\mathrm{A}}}.\tag{21}
\]

我们有 \(\psi(\alpha)\psi(\beta)=0\)，因为 \(r^{2}\) 为零；于是由 (19) 与 (20) 我们推出

\[
(\sigma + \psi) (\alpha \beta) = (\sigma + \psi) (\alpha) (\sigma + \psi) (\beta),\tag{22}
\]

于是 K-线性截面 \(\sigma +\psi\)（\(\pi\) 的）属于 \(\Sigma\)。它的像是一个子代数 \(S\)（\(A\) 的），满足 \(A = S + \mathfrak{r}\)。

我们现在来证明一般情形中 S 的存在性。我们对最小整数 \(p \geqslant 1\)（使 \(r^{p} = 0\)）作归纳推理；p = 1 的情形是平凡的。假设 \(p \geqslant 2\)，并令 \(A' = A/r^{p-1}\)。记 \(r'\) 为 \(A'\) 的根，它等于 \(r/r^{p-1}\)（VIII, 第 155 页的命题 5），故满足 \(r'^{p-1} = 0\)，且代数 \(A'/r'\) 同构于 \(\overline{A} = A/r\)，从而是绝对半单的。由归纳假设，存在子代数 \(S'\)（\(A'\) 的），使得 \(A' = S' \oplus r'\)。于是 \(S'\) 具有形式 \(A''/r^{p-1}\)，其中 \(A''\) 是 A 的一个包含 \(r^{p-1}\) 的子代数，且我们有

\[
\mathrm{A} = \mathrm{A} ^ {\prime \prime} + \mathfrak {r}, \quad \mathfrak {r} ^ {p - 1} = \mathrm{A} ^ {\prime \prime} \cap \mathfrak {r}.\tag{23}
\]

代数 \(A^{\prime\prime}/r^{p-1}\) 同构于 \(A^{\prime}/r^{\prime}\)；我们有 \((\mathfrak{r}^{p-1})^{2}=0\)，故 \(r^{p-1}\) 是 \(A^{\prime\prime}\) 的根。由刚才处理过的情形，存在 \(A^{\prime\prime}\) 的一个子代数 S，使得 \(A^{\prime\prime}=S\oplus r^{p-1}\)；由 (23) 我们推出关系 \(A=S\oplus r\)。

推论 1（韦德伯恩定理）. —— 设 K 是一个交换域，A 是一个 K-代数，r 是 A 的根。假设 K-代数 A/r 是绝对半单的。

\begin{enumerate}
\def\labelenumi{\alph{enumi})}
\item
  设 \(S_1\) 与 \(S_2\) 是 A 的满足 \(A = S_1 \oplus \mathfrak{r} = S_2 \oplus \mathfrak{r}\) 的子代数。存在 A 的一个把 \(S_1\) 变到 \(S_2\) 的特殊自同构。
\item
  若 \(\mathfrak{r}\) 是幂零的，则存在一个子代数 \(S\)（\(A\) 的），满足 \(A = S \oplus \mathfrak{r}\)。
\end{enumerate}

这由命题 7 与 VIII, 第 238 页的定理 2 得出。

推论 2. —— 设 A 是完满域 K 上的一个交换的有限次数代数，r 是它的根。存在 A 的唯一的子代数 S，使得 \(A = S \oplus r\)。此外，S 同构于 K 的有限多个有限次数扩张的一个乘积。

K-代数 A/r 是半单的（VIII, 第 173 页，命题 1）且次数有限；由于域 K 是完满的，它是绝对半单的（VIII, 第 232 页，定理 1）。由于理想 r 是幂零的且 A 是交换的，S 的存在性与唯一性便由推论 1 得出。由于 S 是半单、交换且次数有限的，最后的断言是 VIII, 第 137 页的命题 3 的一个推论。

注. —— 1) A/r 绝对半单这一假设在推论 1 中是本质的（VIII, 第 246 页，习题 4）。

\begin{enumerate}
\def\labelenumi{\arabic{enumi})}
\setcounter{enumi}{1}
\tightlist
\item
  假设 A 是域 K 上的一个阿廷代数。若 A 是交换的，则我们可以证明（VIII, 第 180 页，习题 9）A 同构于一个代数乘积 \(A_{1} \times \cdots \times A_{n}\)，使得对每个 i，\(\mathrm{A}_{i}/\Re(\mathrm{A}_{i})\) 都是一个域。反之，若 A 不是交换的，则 A 可能不同构于这样的代数乘积 \(A_{1} \times \cdots \times A_{n}\)：对每个 i，\(\mathrm{A}_{i}/\Re(\mathrm{A}_{i})\) 都是单环（VIII, 第 247 页，习题 5）。
\end{enumerate}

\BourbakiExercises{习题}

\begin{enumerate}
\def\labelenumi{\arabic{enumi})}
\tightlist
\item
  设 \(\mathrm{K}\) 是一个交换域，\(\mathrm{A}\) 是一个 K-代数。若对每个扩张 \(\mathrm{L}\)（\(\mathrm{K}\) 的），\(\mathrm{A}_{(\mathrm{L})}\) -模 \(\mathrm{M}_{(\mathrm{L})}\) 都无根，则称 A-模 M 是绝对无根的。我们说代数 \(\mathrm{A}\) 是绝对无根的，如果 A-模 \(\mathrm{A}_s\) 是绝对无根的。
\end{enumerate}

\begin{enumerate}
\def\labelenumi{\alph{enumi})}
\item
  一个绝对无根的交换代数是可分的（V, §15, No.~2, 第 119 页, 定义 1）。
\item
  绝对无根的模的每个子模都是绝对无根的；绝对无根的模的一个直和是绝对无根的。
\item
  设 M 是一个有限生成的 A-模；则 M 绝对无根当且仅当存在一个完满域 P（K 的代数扩张），使得 \(\mathfrak{R}_{\mathrm{A}_{(\mathrm{P})}}(\mathrm{M}_{(\mathrm{P})}) = 0\)。
\item
  半单 A-模 M 绝对无根当且仅当属于 M 的支集的每个单模都绝对无根。
\item
  设 S 是一个单 A-模，D 是它的交换子。证明下述性质等价：
\end{enumerate}

\begin{enumerate}
\def\labelenumi{(\roman{enumi})}
\item
  A-模 S 是绝对无根的。
\item
  K-代数 \(\mathrm{D}\) 是绝对无根的。
\item
  \(\mathrm{D}\) 的中心是 \(\mathrm{K}\) 的一个可分扩张。
\end{enumerate}

\(f\) ) 一个半单代数绝对无根当且仅当它的每个单分量的中心都是 K 的一个可分扩张。

\(g\) ) 证明一个绝对无根且在 \(\mathbf{K}\) 上有限维的 A-模是绝对半单的。

\begin{enumerate}
\def\labelenumi{\arabic{enumi})}
\setcounter{enumi}{1}
\tightlist
\item
  设 \(A_{1}\) 与 \(A_{2}\) 是 K-代数，A 是代数 \(A_{1} \otimes_{K} A_{2}\)，\(M_{i}\) 是一个 \(A_{i}\) -模（对 i = 1, 2），M 是 A-模 \(M_{1} \otimes_{K} M_{2}\)。
\end{enumerate}

\begin{enumerate}
\def\labelenumi{\alph{enumi})}
\item
  假设 \(A_{1}\) -模 \(M_{1}\) 绝对无根。证明包含关系 \(\mathfrak{R}_{\mathrm{A}}(\mathrm{M}) \subset \mathrm{M}_{1} \otimes_{\mathrm{K}} \mathfrak{R}_{\mathrm{A}_{2}}(\mathrm{M}_{2})\)。
\item
  若模 \(M_{1}\) 与 \(M_{2}\) 都绝对无根，则 A-模 M 亦然。
\end{enumerate}

\begin{enumerate}
\def\labelenumi{\arabic{enumi})}
\setcounter{enumi}{2}
\item
  设 K 是一个交换域，A 与 B 是单 K-代数。假设 A 在其中心 Z 上秩有限，Z 是 K 的一个代数扩张，且 K-代数 B 绝对无根（习题 1）。证明环 \(A \otimes_{K} B\) 是绝对平坦的（VIII, 第 171 页, 习题 27；注意 \(A \otimes_{K} B\) 的每个元素都属于形如 \(A_{1} \otimes_{K} B\) 的子代数，其中 \(A_{1}\) 是 K 上秩有限的单代数）。
\item
  设 \(\mathrm{K}\) 是一个特征为 \(p > 0\) 的交换域，\(a\) 是 \(\mathrm{K} - \mathrm{K}^p\) 的一个元素。用 A 表示（有限维）K-代数 \(\mathrm{K}[\mathrm{X}] / ((\mathrm{X}^p - a)^2)\)。证明 A 的根 \(\mathfrak{r}\) 的平方为零，K-代数 \(\mathrm{A} / \mathfrak{r}\) 同构于 K 的根式扩张 \(\mathrm{L} = \mathrm{K}[\mathrm{X}] / (\mathrm{X}^{p} - a)\)，且 A 不包含任何同构于 L 的 K-子代数。
\item
  设 \(\mathrm{K}\) 是一个交换域，\(\mathrm{L}\) 是一个有限集。用 B 表示 VIII, 第 169 页习题 17 中构造的代数，把中心元素 \(1 - \sum_{\lambda} c_{\lambda} + \sum_{\lambda \neq \mu} r_{\lambda \mu}\)（\(\mathrm{B}\) 的）记作 \(\varepsilon\)，并用 A 表示代数 \(\mathrm{B} / \mathrm{B}\varepsilon\)。K-代数 A 是有限维的，它的根 \(\mathfrak{r}\) 的平方为零，且 \(\mathrm{A} / \mathfrak{r}\) 同构于 \(\mathrm{K}^{\mathrm{L}}\)。证明 A 不同构于两个非零代数的乘积。
\end{enumerate}

¶ 6) 设 K 是一个交换域，E 是 K 的一个扩张。证明下述性质等价：

\begin{enumerate}
\def\labelenumi{(\roman{enumi})}
\item
  对 K 的每个可分扩张 L，环 \(E_{(L)}\) 是半单的。
\item
  域 \(\mathrm{E}\) 是 \(\mathrm{K}\) 的一个有限次数可分扩张的纯不可分扩张。
\end{enumerate}

（假设 (i)，利用 V, §17, No.~2, 第 140 页的推论与 V, §2, 第 146 页的习题 4 证明 E 是 K 的一个代数扩张。然后证明 K 在 E 中的可分闭包 \(E_{s}\) 次数有限。为此，注意在相反的情形，环 \(\mathrm{E}_{(\mathrm{L})}\)（其中 L 是 K 的一个可分闭包）将包含 n 任意大的正交幂等元族 \((e_{1},\ldots,e_{n})\)。假设 (ii)，证明若 \(\mathrm{E}_{s(\mathrm{L})}\) 同构于 \(\prod_{j=1}^{r}C_{j}\)，则 \(\mathrm{E}_{(\mathrm{L})}\) 同构于 \(\prod_{j=1}^{r}D_{j}\)，其中 \(D_{j}\) 是 \(C_{j}\) 的一个纯不可分扩张。）

¶ 7) 设 E 是特征 p \textgreater{} 0 的交换域 K 的一个 p-根式扩张。假设环 \(E \otimes_{K} E\) 是诺特的。

\begin{enumerate}
\def\labelenumi{\alph{enumi})}
\item
  证明存在一个整数 \(k\)，使得 \(\mathrm{E} \subset \mathrm{K}^{p^{-k}}\)（用反证法推理：构造一个序列 \((x_{n})\)（在 \(\mathrm{E}\) 中），使得 \((x_{n} \otimes 1 - 1 \otimes x_{n})^{p^{n-1}} \neq 0\) 且 \((x_{n} \otimes 1 - 1 \otimes x_{n})^{p^n} = 0\)）。
\item
  设 \(\mathrm{F} = \mathrm{K}(\mathrm{E}^{p})\)，并设 \((x_{i})_{i \in \mathrm{I}}\) 是 E 在 F 上的一个 p-基（V, §13, No.~1, 第 98 页）。令 \(y_{i} = x_{i} \otimes 1\) 与 \(z_{i} = x_{i} \otimes 1 - 1 \otimes x_{i}\)。设 \(\Lambda\) 是 \(\mathbf{N}^{(I)}\) 的由族 \((\alpha_{i})_{i \in \mathrm{I}}\)（满足 \(\alpha_{i} < p\)，对每个 i）组成的子集。证明元素 \(\prod_{i \in I} y_{i}^{\alpha_{i}} z_{i}^{\beta_{i}}\)（对 \((\alpha_{i})\) , \((\beta_{i}) \in \Lambda\)）构成 F-向量空间 \(E \otimes_{F} E\) 的一个基。由此得出 I 是有限的（注意环 \(E \otimes_{F} E\) 是诺特的）。
\item
  由此得出 E 在 K 上次数有限（利用 V, §13, 第 170 页, 习题 1, b)）。
\end{enumerate}

¶ 8) 设 E 是交换域 K 的一个代数扩张。证明环 \(E \otimes_{K} E\) 是诺特的当且仅当 E 在 K 上次数有限（利用习题 6 的方法化归为纯不可分扩张的情形，然后应用习题 7）。

¶ 9) 设 K 是一个交换域，A 是一个 K-代数。假设对每个阿廷 K-代数 B，环 \(A \otimes_{K} B\) 都是阿廷的。证明 A 在 K 上有限维（通过注意 \(A/\Re(A)\) -模 \(\Re(A)^{k}/\Re(A)^{k+1}\) 有有限长度，化归为 A 半单的情形；然后利用习题 6 与 7）。

\begin{enumerate}
\def\labelenumi{\arabic{enumi})}
\setcounter{enumi}{9}
\tightlist
\item
  设 K 是一个交换环，A 是一个 K-代数，P 是一个 (A,A)-双模。A 的一个经 P 的扩张是一个三元组 \((B,i,p)\)，其中 B 是一个 K-代数，\(p:B\to A\) 是 K-代数的一个满同态，而 \(i:P\to B\) 是象为 Ker p 的单 K-线性同态，满足 \(i(p(b)x)=bi(x)\) 与 \(i(xp(b))=i(x)b\)（对 \(x\in P\) 与 \(b\in B\)）。从扩张 \((B,i,p)\) 到扩张 \((B',i',p')\) 的一个同构是 K-代数的一个同构 \(u:B\to B'\)，使得 \(p'\circ u=p\) 且 \(u\circ i=i'\)。
\end{enumerate}

\begin{enumerate}
\def\labelenumi{\alph{enumi})}
\item
  设 \((B, i, p)\) 是 A 的一个经 P 的扩张；i 的象是 B 的一个平方为零的双边理想。反之，若 \(p : B \to A\) 是 K-代数的一个满同态且其核 P 的平方为零，而 i 是 P 到 B 中的典范单射，则三元组 \((B, i, p)\) 是 A 的一个经 P 的扩张。
\item
  设 \(i_{0}: P \to A \oplus P\) 与 \(p_{0}: A \oplus P \to A\) 是典范映射。在 \(A \oplus P\) 上存在唯一的 K-代数结构，使得 \((A \oplus P, i_{0}, p_{0})\) 是 A 的一个经 P 的扩张。证明这个扩张的自同构群可以等同于从 A 到 P 的 K-导子所成的群。
\item
  设 \((B, i, p)\) 是 A 的一个经 P 的扩张。证明下述条件等价：
\end{enumerate}

\begin{enumerate}
\def\labelenumi{(\roman{enumi})}
\item
  扩张 \((\mathrm{B},i,p)\) 同构于扩张 \((\mathrm{A}\oplus \mathrm{P},i_0,p_0)\)。
\item
  存在 K-代数的一个同态 \(s: \mathrm{A} \to \mathrm{B}\)，使得 \(p \circ s = \operatorname{Id}_{\mathrm{A}}\)。
\item
  存在子代数 \(S\)（\(B\) 的），使得 \(B = S \oplus i(P)\)。
\end{enumerate}

若这些条件满足，则我们说扩张 \((B,i,p)\) 是平凡的。d) 设 \((B,i,p)\) 与 \((B',i',p')\) 是 A 的两个经 P 的扩张。设 C 是 \(B \times B'\) 的由元素 \((b,b')\)（满足 \(p(b) = p'(b')\)）组成的子代数，设 c 是 C 的由元素 \((i(x), -i'(x))\)（对 \(x \in P\)）组成的双边理想，并设 \(B''\) 是代数 C/c 且 \(\pi : C \to B''\) 是典范同态。设 \(i'' : P \to B''\) 与 \(p'' : B'' \to A\) 是由 \(i''(x) = \pi(i(x),0)\) 与 \(p''(\pi(b,b')) = p(b)\)（对 \(x \in P\) , \(b \in B\) , \(b' \in B'\)）定义的同态。证明 \((B'', i'', p'')\) 是 A 的一个经 P 的扩张；它称为 \((B,i,p)\) 与 \((B',i',p')\) 的融合和。证明我们这样就在 A 的经 P 的扩张的同构类所成的集合 \(\mathrm{Ex}_{\mathrm{K}}(\mathrm{A},\mathrm{P})\) 上定义了一个交换群律，其单位元是平凡扩张的类。此外，K 在 \(\mathrm{Ex}_{\mathrm{K}}(\mathrm{A},\mathrm{P})\) 上由 \(\lambda(\mathrm{B},i,p) = (\mathrm{B},\lambda^{-1}i,p)\)（对 \(\lambda \in K\) , \(\lambda \neq 0\)）给出的作用在 \(\mathrm{Ex}_{\mathrm{K}}(\mathrm{A},\mathrm{P})\) 上定义了一个 K-模结构。

\begin{enumerate}
\def\labelenumi{\alph{enumi})}
\setcounter{enumi}{4}
\tightlist
\item
  设 \(u: P \to P'\) 是 (A, A)-双模的一个同态，并设 \((B, i, p)\) 是 A 的一个经 P 的扩张。设 \(B'\) 是融合和 \(B \oplus_{P} P'\)（II, §1, 第 381 页, 习题 5）；对 \(b \in B\) , \(x' \in P'\)，把 \((b, x')\) 在 \(B'\) 中的类记作 \([b, x']\)。设 \(i': P' \to B'\) , \(p': B' \to A\) 与 \(v: B \to B'\) 是由下式定义的映射
\end{enumerate}

\[
i ^ {\prime} (x ^ {\prime}) = [ 0, x ^ {\prime} ], \qquad p ^ {\prime} ([ b, x ^ {\prime} ]) = p (b), \qquad v (b) = [ b, 0 ].
\]

证明在 \(B'\) 上存在唯一的代数结构，使得 \((\mathrm{B}', i', p')\) 是 A 的一个经 \(P'\) 的扩张且 v 是代数同态。我们说 \((\mathrm{B}', i', p')\) 是 A 的经 \(P'\) 的、由 \((\mathrm{B}, i, p)\) 经 u 得出的扩张；每个扩张 \((\mathrm{C}, j, q)\)（A 的经 \(P'\) 的扩张），若存在代数同态 \(w: B \to C\) 满足 \(q \circ w = p\) 与 \(w \circ i = j \circ v\)，则它同构于 \((\mathrm{B}', i', p')\)。

\(f\) ) 映射 \(\mathrm{Hom}_{(\mathrm{A},\mathrm{A})}(\mathrm{P},\mathrm{P}')\times \mathrm{Ex}_\mathrm{K}(\mathrm{A},\mathrm{P})\to \mathrm{Ex}_\mathrm{K}(\mathrm{A},\mathrm{P}')\) 把二元组 \((u,(B,i,p))\) 映到由 (B, \(i,p)\) 经 \(u\) 得出的扩张的类，它是 K-双线性的。\(g\) ) 设 \(f:\mathrm{A}'\to \mathrm{A}\) 是一个 K-代数同态，并设 (B, \(i,p)\) 是 A 的一个经 P 的扩张。类似地定义纤维积 \(\mathrm{B}'' = \mathrm{B}\times_{\mathrm{A}}\mathrm{A}'\) 上的代数结构以及扩张 \((\mathrm{B}'' , i'', p'')\)（\(\mathrm{A}'\) 的经 P 的扩张），它称为由 (B, \(i,p)\) 经 \(f\) 得出的扩张。

\begin{enumerate}
\def\labelenumi{\alph{enumi})}
\setcounter{enumi}{7}
\tightlist
\item
  再假设代数 A 是交换的。证明 A 的经 P 的、使得代数 B 为交换的扩张 \((B,i,p)\) 的类构成 \(\operatorname{Ex}_{\mathrm{K}}(\mathrm{A},\mathrm{P})\) 的一个 K-子模，我们把它记作 \(\operatorname{Ex}_{\mathrm{K}}^{c}(\mathrm{A},\mathrm{P})\)。
\end{enumerate}

\begin{enumerate}
\def\labelenumi{\arabic{enumi})}
\setcounter{enumi}{10}
\tightlist
\item
  \begin{enumerate}
  \def\labelenumii{\alph{enumii})}
  \tightlist
  \item
    设 f 是 \(\mathrm{C}^{2}(\mathrm{A},\mathrm{P})\) 的一个元素（VIII, 第 239 页）。证明我们在 K-模 \(A \oplus P\) 上通过令 \((a,x)(a',x') = (aa',ax' + xa' + f(a,a'))\) 来定义一个代数结构。典范单射 \(i: P \to A \oplus P\) 与典范满射 \(p: A \oplus P \to A\) 是环同态，且三元组 \((\mathrm{A} \oplus \mathrm{P}, i, p)\) 是 A 的一个经 P 的扩张（注意我们有 \(f(a,1) = af(1,1)\) 与 \(f(1,a) = f(1,1)a\)（对每个 \(a \in A\)））。
  \end{enumerate}
\end{enumerate}

\begin{enumerate}
\def\labelenumi{\alph{enumi})}
\setcounter{enumi}{1}
\item
  证明这个扩张的同构类只依赖于 f 在 \(\mathrm{H}^{2}(\mathrm{A},\mathrm{P})\) 中的类。于是 a) 中的构造定义了一个从 \(\mathrm{H}^{2}(\mathrm{A},\mathrm{P})\) 到 \(\mathrm{Ex}_{\mathrm{K}}(\mathrm{A},\mathrm{P})\) 的由作为 K-模的扩张为平凡的 A 的经 P 的扩张组成的 K-子模上的 K-模同构。
\item
  假设环 A 是交换的；用 \(\mathrm{C}^{2}(\mathrm{A},\mathrm{P})^{s}\) 表示 \(\mathrm{C}^{2}(\mathrm{A},\mathrm{P})\) 的由满足 \(f(x,y)=f(y,x)\)（对 A 中所有 x, y）的元素 f 组成的 K-子模，并用 \(\mathrm{H}^{2}(\mathrm{A},\mathrm{P})^{s}\) 表示 \(\mathrm{C}^{2}(\mathrm{A},\mathrm{P})^{s}\) 在 \(\mathrm{H}^{2}(\mathrm{A},\mathrm{P})\) 中的象。证明 b) 中的同构诱导一个从 \(\mathrm{H}^{2}(\mathrm{A},\mathrm{P})^{s}\) 到 \(\mathrm{Ex}_{\mathrm{K}}^{c}(\mathrm{A},\mathrm{P})\) 的由作为 K-模的扩张为平凡的交换的 A 的经 P 的扩张组成的 K-子模的同构。
\end{enumerate}

\begin{enumerate}
\def\labelenumi{\arabic{enumi})}
\setcounter{enumi}{11}
\tightlist
\item
  设 A 是一个交换环，J 是 A 的一个理想，B 是环 A/J，P 是一个 B-模。a) 证明正合列 \(0 \rightarrow J/J^{2} \rightarrow A/J^{2} \rightarrow B \rightarrow 0\) 定义了 A-代数 B 的一个经 \(J/J^{2}\) 的扩张（习题 10）。
\end{enumerate}

\begin{enumerate}
\def\labelenumi{\alph{enumi})}
\setcounter{enumi}{1}
\tightlist
\item
  设 P 是一个 B-模。证明把 u 映到由 a) 中的扩张经 u 得出的扩张的同态 \(\operatorname{Hom}_{\mathrm{B}}(\mathrm{J}/\mathrm{J}^{2},\mathrm{P}) \to \operatorname{Ex}_{\mathrm{A}}^{c}(\mathrm{B},\mathrm{P})\) 是一个同构。
\end{enumerate}

\section{中心单代数}

在本节中，K 是一个交换域。

\subsection{中心单代数}

定义 1. —— 若映射 \(\lambda\mapsto\lambda1\) 是从 K 到 A 的中心的一个双射，则称域 K 上的代数 A 是中心的。

中心代数不化为 0。对每个整数 \(n \geqslant 1\)，矩阵 K-代数 \(\mathbf{M}_{n}(\mathrm{K})\) 是中心的（VIII, 第 83 页, 推论 2）且是单的（VIII, 第 120 页, 定理 1）。更一般地，设 D 是一个有限次数的中心 K-代数；则 \(\mathbf{M}_{n}(\mathrm{D})\) 也是中心的。设 A 是一个单环；它的中心 Z 是一个域（VIII, 第 121 页, 推论 1），因此 A 是 Z 上的一个中心单代数。若域 K 是代数闭的，则 K 上的一个有限次数的单代数是中心的（VIII, 第 122 页, 推论 3）。中心单代数的反代数是中心单的。

注. —— 1) 设 A 与 B 是 K-代数。若 \(A \otimes_{K} B\) 是一个中心单代数，则 A 与 B 亦然（III, §4, No.~4, 第 468 页, 命题 6 的推论与 VIII, 第 221 页, 命题 6）。反之，若 A 与 B 都是中心单代数且其中之一在 K 上次数有限，则 \(A \otimes_{K} B\) 是一个中心单代数（III, §4, No.~4, 第 468 页, 命题 6 的推论与 VIII, 第 222 页, 推论 2）。

\begin{enumerate}
\def\labelenumi{\arabic{enumi})}
\setcounter{enumi}{1}
\item
  设 A 是一个 K-代数，L 是域 K 的一个扩张。若 L-代数 \(A_{(L)}\) 是中心单的，则 K-代数 A 是中心单的。反之，若次数 {[}A : K{]} 与 {[}L : K{]} 之一有限且 A 是一个中心单 K-代数，则 L-代数 \(A_{(L)}\) 是中心单的。这由 VIII, 第 222 页的推论 2 得出。
\item
  设 A 与 B 是森田等价的 K-代数。代数 A 是中心单的当且仅当 B 是中心单的（VIII, 第 105 页, 命题 6；第 111 页, 推论；以及第 102 页, 推论 1）。
\item
  特别地，若 A 是一个中心单 K-代数且 \(n \geqslant 1\)，则 \(\mathbf{M}_{n}(\mathrm{A})\) 是一个中心单 K-代数（VIII, 第 102 页, 例 1）。
\end{enumerate}

定理 1. —— 设 A 是一个有限次数的 K-代数。下述性质等价：

\begin{enumerate}
\def\labelenumi{(\roman{enumi})}
\item
  代数 A 是中心单的。
\item
  代数 \(\mathbf{A}\) 是中心的且无根。
\item
  从 K-代数 \(A \otimes_{K} A^{\circ}\) 到 K-代数 \(\operatorname{End}_{K}(A)\) 的把 \(a \otimes a'\) 映到从 A 到 A 的 K-线性映射 \(x \mapsto axa'\) 的典范同态是双射的。
\item
  存在域 K 的一个扩张 L 与一个整数 \(n \geqslant 1\)，使得 L-代数 \(\mathrm{A}_{(\mathrm{L})}\) 与 \(\mathbf{M}_{n}(\mathrm{L})\) 同构。
\item
  对每个可分闭包 \(\mathrm{K}'\)（\(\mathbf{K}\) 的），存在一个整数 \(n \geqslant 1\)，使得 \(\mathrm{K}'\) -代数 \(\mathrm{A}_{(\mathrm{K}')}\) 与 \(\mathbf{M}_n(\mathrm{K}')\) 同构。
\item
  存在域 K 的一个有限次数的伽罗瓦扩张 L 与一个整数 \(n \geqslant 1\)，使得 L-代数 \(\mathrm{A}_{(\mathrm{L})}\) 与 \(\mathbf{M}_n(\mathrm{L})\) 同构。
\item
  存在一个有限次数的 K-代数 D（它是一个以 K 为中心的域）与一个整数 \(n \geqslant 1\)，使得代数 A 同构于代数 \(\mathbf{M}_{n}(\mathrm{D})\)。
\end{enumerate}

一个环是单的当且仅当它是半单的且它的中心是一个域（VIII, 第 143 页, 命题 10 的推论）。由于 A 是域 K 上的一个有限次数代数，它是一个左阿廷环；因此它半单当且仅当它无根（VIII, 第 154 页, 命题 4）。由此得出 (i) 与 (ii) 的等价性。

令 \(E = A \otimes_{K} A^{\circ}\) 与 \(F = \text{End}_{K}(A)\)；用 \(\varphi\) 表示由关系式 \(\varphi(a \otimes a')(x) = axa'\)（对 A 中的 \(x, a, a'\)）定义的从 E 到 F 的典范同态。若代数 A 是中心单的，则 \(A^{\circ}\) 亦然，从而 E 亦然（注 1），于是 \(\varphi\) 是单射的。现在我们有 \([E : K] = [A : K]^{2} = [F : K]\)，故 \(\varphi\) 是双射的。反之，假设 \(\varphi\) 是双射的；由于代数 F 是中心单的（因为它同构于一个矩阵代数 \(\mathbf{M}_{m}(K)\)），E 亦然，从而 A 亦然（注 1）。我们这样就证明了 (i) 与 (iii) 的等价性。

由注 4，断言 (vii) 蕴涵断言 (i)。其逆蕴涵由 VIII, 第 122 页的推论 3 与 VIII, 第 83 页的推论 2 得出。

显然 (vi) 蕴涵 (iv)，而由注 2，(iv) 蕴涵 (i)。

剩下要证明蕴涵关系 (i)⇒(v)⇒(vi)。假设 A 是中心单的，并设 \(K'\) 是 K 的一个可分闭包（V, §7, No.~8, 第 45 页）。则 \(\mathrm{A}_{(\mathrm{K}')}\) 是 \(K'\) 上的一个有限次数的中心单代数（VIII, 第 251 页, 注 2）。于是由 VIII, 第 231 页的推论，存在一个整数 \(n \geqslant 1\) 与一个 \(K'\) -代数同构（从 \(\mathrm{A}_{(\mathrm{K}')}\) 到 \(\mathbf{M}_{n}(\mathrm{K}')\)）；注意 \(K'\) -代数 \(\mathbf{M}_{n}(\mathrm{K}')\) 与 \(\mathbf{M}_{n}(\mathrm{K})_{(\mathrm{K}')}\) 同构。由 VIII, 第 234 页的引理 4，存在 \(K'\) 的一个在 K 上有限生成的子扩张 L，使得 L-代数 \(\mathrm{A}_{(\mathrm{L})}\) 与 \(\mathbf{M}_{n}(\mathrm{K})_{(\mathrm{L})}\) 同构。于是 L 在 K 上是可分的且次数有限，故包含在一个在 K 上伽罗瓦且次数有限的子扩张 \(L'\)（\(K'\) 的）中（V, §10, No.~1, 第 57 页, 命题 2）。于是 \(L'\) -代数 \(\mathrm{A}_{(\mathrm{L}')}\) 与 \(\mathbf{M}_{n}(\mathrm{L}')\) 同构。

推论 1. —— 设 A 是可分闭域 K 上的一个有限次数的中心单代数。存在一个整数 \(n \geqslant 1\)，使得 A 同构于矩阵代数 \(\mathbf{M}_{n}(\mathrm{K})\)。

实际上，\(\mathrm{K}\) 的每个伽罗瓦扩张都等于 \(\mathrm{K}\)；只需应用定理 1 的性质 (i) 与 (v) 的等价性。

推论 2. —— 设 A 是 K 上的一个有限次数的中心单代数（例如，K 上的一个以 K 为中心的有限次数的域）。存在一个整数 \(n \geqslant 1\)，使得 \([A : K] = n^{2}\)。

设 \(\mathrm{L}\) 是 \(\mathrm{K}\) 的一个扩张，\(n\) 是一个严格正的整数，使得 L-代数 \(\mathrm{A}_{(\mathrm{L})}\) 与 \(\mathbf{M}_n(\mathrm{L})\) 同构。我们有

\[
[ \mathrm{A}: \mathrm{K} ] = [ \mathrm{A} _ {(\mathrm{L})}: \mathrm{L} ] = [ \mathbf {M} _ {n} (\mathrm{L}): \mathrm{L} ] = n ^ {2}.
\]

在推论 2 的记号中，整数 n 称为 A 的约化次数。

注 5. —— 设 A 是 K 上的一个约化次数为素数 \(\ell\) 的有限次数中心单代数。则 A 或者是一个域，或者同构于 \(\mathbf{M}_{\ell}(\mathrm{K})\)。实际上，A 同构于形如 \(\mathbf{M}_{n}(\mathrm{D})\) 的一个代数，其中 D 是一个以 K 为中心的域，且我们有

\[
\ell^ {2} = [ \mathrm{A}: \mathrm{K} ] = n ^ {2} [ \mathrm{D}: \mathrm{K} ];
\]

若 A 不是一个域，则 \(n \neq 1\)，故 \(n = \ell\) 且 D = K。

\subsection{关于双模的两个引理}

设 A 与 B 是环。对从 B 到 A 的任一同态 f，我们用 \(A^{f}\) 表示以右 A-模 \(A_{d}\) 为基础、其左 B-模结构的外部法则由 \((b,a)\mapsto f(b)a\) 给出的 (B,A)-双模。

引理 1. —— 设 f 与 g 是从 B 到 A 的同态。下述条件等价：

\begin{enumerate}
\def\labelenumi{(\roman{enumi})}
\item
  (B,A)-双模 \(A^{f}\) 与 \(A^{g}\) 同构。
\item
  存在 A 的一个内自同构 \(\theta\)（I, §8, No.~4, 第 102 页, 例 2），使得 \(g = \theta \circ f\)。
\end{enumerate}

右 A-模 \(A_{d}\) 的自同构是映射 \(x \mapsto ax\)，其中 a 是 A 的一个可逆元。这样的一个自同构是从 \(A^{f}\) 到 \(A^{g}\) 的 B-线性映射当且仅当我们有

\[
g (b) a x = a f (b) x
\]

对 A 中的每个 x 与 B 中的每个 b。这个关系式等价于对 B 中的每个 b 有 \(g(b) = af(b)a^{-1}\)，即等价于 \(g = \theta \circ f\)，其中 \(\theta\) 是 A 的内自同构 \(x \mapsto axa^{-1}\)。

引理 2. —— 假设 B 是一个半单环，且是其中心 Z 上的一个有限生成模。设 M 与 N 是 (B, A)-双模。假设它们有有限长度（特别地，当它们是有限长度的右 A-模时就是这种情形）。若 M 与 N 作为 (Z, A)-双模同构，则它们作为 (B, A)-双模同构。

\begin{enumerate}
\def\labelenumi{\Alph{enumi})}
\item
  首先考虑 B 是交换域 L 上一个有限维 d 向量空间 S 的自同态环的情形。这时 Z = L；我们把 S 看作一个 (B, Z)-双模。环 B 是单的，S 是一个单 B-模，而 Z 是 S 的交换子；每个 B-模都是 S 型的同型模（VIII, 第 122 页, 命题 2 a)）。设 \((V, \alpha)\)（相应地 \((W, \beta)\)）是 B-模 M（相应地 N）的一个描述。集合 V（相应地 W）带有一个 (Z, A)-双模结构，使得 \(\alpha\)（相应地 \(\beta\)）是 (B, A)-双模的同构（VIII, 第 64 页, 注 2）。作为 (Z, A)-双模，M 同构于 \(V^{d}\) 而 N 同构于 \(W^{d}\)，且存在一个从 V 的 (Z, A)-子双模的集合（按包含排序）到 M 的 (B, A)-子双模的集合的同构（同前所引）。于是 V 是一个有限长度的 (Z, A)-双模，W 亦然。由于 (Z, A)-双模 \(V^{d}\) 与 \(W^{d}\) 同构，由应用于环 \(Z \otimes_{Z} A^{\circ}\) 的 VIII, 第 37 页的定理 2, d)，(Z, A)-双模 V 与 W 同构。最后，(B, A)-双模 M 与 N 同构。
\item
  现在考虑 B 是一个作为 Z-模有限生成的单环的情形。这时 Z 是一个域，而 B 是域 Z 上的一个有限次数的中心单代数。由 VIII, 第 252 页的定理 1，存在 Z 的一个在 Z 上次数有限的扩张 \(Z'\)，使得 \(Z'\) -代数 \(B' = B_{(Z')}\) 同构于一个有限维 \(Z'\) -向量空间的自同态代数。令 \(M' = M_{(Z')}\) 与 \(N' = N_{(Z')}\)。则 \(M'\) 与 \(N'\) 是有限长度的 \((B', A)\) -双模；看作 \((Z', A)\) -双模时，\(M'\) 与 \(N'\) 同构。由 A) 中处理的情形，\(M'\) 与 \(N'\) 作为 \((B', A)\) -双模同构，从而更不用说作为 \((B, A)\) -双模同构。令 \(r = [Z' : Z]\)。\((B, A)\) -双模 \(M' = Z' \otimes_{Z} M\) 同构于 \(M^r\)，同样地，\((B, A)\) -双模 \(N'\) 同构于 \(N^r\)。由于 M 与 N 是有限长度的 \((B, A)\) -双模，由 VIII, 第 37 页的定理 2, d) 可知 \((B, A)\) -双模 M 与 N 同构。
\item
  最后考虑一般情形，即 B 是一个作为 Z-模有限生成的半单环。设 S 是单 B-模的类所成的集合；它是有限的（VIII, 第 136 页, 命题 1）。对任何 \(\lambda \in S\)，用 \(M_{\lambda}\)（相应地 \(N_{\lambda}\)）表示 B-模 M（相应地 N）的 \(\lambda\) 型同型分量；它是 M（相应地 N）的一个 (B, A)-子双模（VIII, 第 67 页的注）。对 \(\lambda \in S\)，把 B-模 \(\lambda\) 的零化子记作 \(b_{\lambda}\)，并令 \(B_{\lambda} = B / b_{\lambda}\)；设 \(Z_{\lambda}\) 是 \(B_{\lambda}\) 的中心。对 \(\lambda \in S\)，\((B_{\lambda}, A)\) -双模 \(M_{\lambda}\) 与 \(N_{\lambda}\) 有有限长度。于是我们可以把 B 等同于单环 \(B_{\lambda}\) 的乘积，把 Z 等同于 \(Z_{\lambda}\) 的乘积（VIII, 第 141 页, 命题 8）。此外，我们可以把 M 等同于 \(\prod_{\lambda \in S} M_{\lambda}\)，把 N 等同于 \(\prod_{\lambda \in S} N_{\lambda}\)。由假设，M 与 N 作为 (Z, A)-双模同构；由此得出，对 \(\lambda \in S\)，\(M_{\lambda}\) 与 \(N_{\lambda}\) 作为 ( \(Z_{\lambda}, A\) )-双模同构。由 B) 中处理的情形，\((B_{\lambda}, A)\) -双模 \(M_{\lambda}\) 与 \(N_{\lambda}\) 同构，于是 (B, A)-双模 M 与 N 同构。
\end{enumerate}

注. —— 由引理 2 的证明可知 M 与 N 是有限长度的 (Z, A)-双模。因此，若 B 与 A 是两个作为其各自的中心 Z(B) 与 Z(A) 上的有限生成模的半单环，则两个作为 (Z(B), Z(A))-双模同构的有限长度 (B, A)-双模是同构的。
